\chapter{信息抽取与知识构建}
\label{chap:nlp-information-extraction}

读出通知中的“4月18日”，还不能直接填写活动日期：原文也可能说“原定4月18日，现已取消”。
找出两个组织名称，也不能据此建立合作关系。上一章讨论的词法、句法、语义和篇章分析帮助
我们辨认语言结构，本章进一步问：怎样把这些结构转成可查询、可核验，并能随新材料更新的记录？

沿用本部分的教学通知：活动编号为\code{ML-0418}，通知发布于2025年4月15日，
时间均为北京时间。编号和发布日期是给定元数据，不是从下列正文中猜出的字段。
\begin{quote}
机器学习读书会原定于4月17日举行，现改至4月18日（周五）14:00，地点为海棠楼201。
校内师生可直接参加，校外来宾须在4月16日18:00前提交报名信息。
\end{quote}
本章所有通知、知识库和数值例子均为教学构造，不表示真实活动或模型实测结果。
需要额外的组织、人物或业务状态时，会在当地明确给出补充材料。

\term{信息抽取}（information extraction, IE）按照给定的类型、角色和字段规范，
把原文表达的对象与事实恢复为结构化结果。规范通常称为\term{模式}（schema）；
规定对象类型、关系含义和角色约束的概念体系称为\term{本体}（ontology）。
模式可以要求输出地点，却不能替原文补出地点；模型可以判断“改至”指向新日期，
却还需留下它依据哪个句子、哪个位置作出判断。

本章从不同记录对象分解这个问题。实体抽取确定提及位置和对象身份，关系抽取确定对象之间
说了什么，事件抽取把参与者和条件绑定到同一次事件，事件关系描述各次事件的时间、因果与
组成联系。知识记录的构建则检查这些输出能否对齐、合并和更新。几项任务可以联合预测，
但验收对象不随网络是否共享而消失。表~\ref{tab:nlp-extraction-errors}给出错误定位的入口：
同一条记录可能同时有位置、身份和条件错误，不宜只保留一个端到端分数。

\begin{table*}[htbp]
  \centering
  \small
  \begin{tabularx}{\linewidth}{@{}>{\raggedright\arraybackslash}p{13mm}
    *{5}{>{\raggedright\arraybackslash}X}
    >{\raggedright\arraybackslash}p{23mm}@{}}
    \toprule
    产出 & 定位错误 & 类型或角色错误 & 身份错配 & 条件遗漏 & 无依据新增 & 验收入口 \\
    \midrule
    实体 & 少取“201” & 地点标为组织 & 不适用 & 丢失修饰语 & 补出主讲人 & 跨度、类型、联合匹配 \\
    链接 & 链错同名位置 & 类型不兼容 & 连到异校同名楼 & 忽略库版本 & 强行取消NIL & 召回、身份、拒绝 \\
    关系 & 证据句不完整 & 举办方向颠倒 & 主客体绑错 & 漏掉否定 & 共现当合作 & 对象、标签、证据 \\
    事件 & 触发或论元错位 & 旧时间当新时间 & 两场活动混合 & 漏掉计划状态 & 填入未给费用 & 提及、角色、完整记录 \\
    事件关系 & 端点事件选错 & 先后当因果 & 误合并端点 & 漏掉参照日期 & 强排未知顺序 & 有向边、证据、一致性 \\
    综合记录 & 来源偏移失效 & 字段映射错误 & 近似记录误并 & 截止时刻缺失 & 冲突值擅自裁决 & 对齐、合并、版本与来源 \\
    \bottomrule
  \end{tabularx}
  \caption{抽取产出与错误诊断。固定同一原文和本体，分别检查各类错误；“不适用”表示
  尚未进行知识库身份选择。局部指标与最终记录指标都需报告，不能用给定候选条件下的
  分类结果代替从原始文本开始的抽取结果。}
  \label{tab:nlp-extraction-errors}
\end{table*}

\section{实体识别与身份确定}
\label{sec:nlp-extraction-entities}

“海棠楼201”有三种需要分开的解释。它在原文中占据一段字符，这属于
\term{实体提及}（entity mention）；补充句中的“该教室”可能与它指向同一对象，
这属于文档内共指；把它对应到校园场所库中的某个编号，才属于实体链接。
提及是语言位置，实体是所指对象，共指簇则是一组被判为同指的提及。
即使共指判断正确，也可能把整簇链接到另一校区的同名教室。

图~\ref{fig:nlp-extraction-identity}使用明确标出的补充通知和小型知识库展示这三项产出。
图中的临时服务点没有库中编号，仍然可以被识别、参与共指和进入带来源的记录；
不应因为尚未登记身份而把它从原文中删除。
\begin{figure*}[htbp]
  \centering
  % Adapted from academic-drawing's TikZ template; project fonts and palette prevail.
% Final canvas: 167 mm wide. Ports: text east -> spans west; spans east
% -> cluster west; cluster east -> KB west. Side branches use a routing lane.
\begin{tikzpicture}[
  x=1mm,y=1mm,
  font=\figurefont\fontsize{9}{11}\selectfont,
  box/.style={rectangle,align=left,inner sep=2mm,line width=0.85pt,
    draw=BookRule,fill=white},
  source/.style={box,text width=31mm,draw=FigureInput,fill=FigureInput!8},
  span/.style={box,text width=27mm,draw=FigureInput,fill=FigureInput!8},
  cluster/.style={box,text width=23mm,draw=FigureModel,fill=FigureModel!8},
  target/.style={box,text width=33mm,draw=FigureOutput,fill=FigureOutput!8},
  locate/.style={draw=FigureInput,line width=0.85pt,rounded corners=1.2mm},
  coref/.style={draw=FigureModel,line width=0.85pt,dashed,rounded corners=1.2mm},
  link/.style={draw=FigureOutput,line width=1.1pt,-{Latex[length=2mm]}},
  head/.style={font=\figurefont\bfseries\fontsize{9}{11}\selectfont,
    text=BookInk,align=center}
]
  \node[head] at (18,37) {补充原文};
  \node[head] at (62,37) {原文跨度};
  \node[head] at (104,37) {文档共指簇};
  \node[head] at (146,37) {教学知识库};
  \node[source] (text) at (18,0) {
    地点为海棠楼201。\\
    该教室旁设临时服务点。\\[3pt]
    字符从0开始计数
  };
  \node[span] (m1) at (62,23) {海棠楼201\\$[3,9)$};
  \node[span] (m2) at (62,0) {该教室\\$[10,13)$};
  \node[span] (m3) at (62,-23) {临时服务点\\$[15,20)$};
  \node[cluster] (c1) at (104,11.5) {同一教室\\$\{m_1,m_2\}$};
  \node[cluster] (c2) at (104,-23) {服务点\\$\{m_3\}$};
  \node[target] (kb1) at (146,11.5) {ROOM-201\\海棠楼201};
  \node[target] (kb2) at (146,-23) {NIL\\无可确认库中项};
  \draw[locate] (text.east) -- (42,0);
  \draw[locate] (42,23) -- (42,-23);
  \draw[locate] (42,23) -- (m1.west);
  \draw[locate] (42,0) -- (m2.west);
  \draw[locate] (42,-23) -- (m3.west);
  \draw[coref] (m1.east) -- (84,23) -- (84,11.5) -- (c1.west);
  \draw[coref] (m2.east) -- (84,0) -- (84,11.5);
  \draw[coref] (m3.east) -- (c2.west);
  \draw[link] (c1.east) -- (kb1.west);
  \draw[link] (c2.east) -- (kb2.west);
  \draw[locate] (4,-38) -- (14,-38);
  \node[anchor=west] at (16,-38) {定位};
  \draw[coref] (49,-38) -- (59,-38);
  \node[anchor=west] at (61,-38) {共指归组};
  \draw[link] (111,-38) -- (121,-38);
  \node[anchor=west] at (123,-38) {身份选择};
\end{tikzpicture}

  \caption{实体提及、共指与身份链接。教学补充材料为“地点为海棠楼201。该教室旁设临时服务点”。
  原文到跨度的实线表示定位，跨度到共指簇的虚线表示同指归组，簇到知识库的箭头表示身份选择。
  临时服务点的NIL表示当前知识库没有可确认的对应项，不表示该对象不存在。}
  \label{fig:nlp-extraction-identity}
\end{figure*}

\FloatBarrier
\subsection{实体提及识别}
\label{sec:nlp-extraction-mentions}

实体识别（named entity recognition, NER）要同时回答“哪里是提及”和“属于什么类型”。
沿用第~\ref{chap:nlp-introduction}章的位置约定，原文字符串的字符区间写成从零开始的
$[s,e)$。模型内部采用从1开始的词元编号，词元跨度$(i,j)$包含两端；
若词元$i$对应原文$[a_i,b_i)$，连续跨度恢复为$[a_i,b_j)$。
子词切分、全半角转换和空格清理都必须保留这一映射，否则模型索引不能直接充当原文偏移。

标注规范决定哪些结构应该保留。补充教学句“海棠大学计算机学院发布通知”可把
“海棠大学”和“海棠大学计算机学院”都标成组织，形成\term{嵌套提及}（nested mention）。
若只允许互不相交的实体，则是平面规范。不连续提及则由多个片段共同表达一个对象，
其输出应是有序片段列表，而不是覆盖所有间隔文字的一个大区间。
图式输出能表示这种关联，普通单层BIO序列不能。

细粒度和少样本是另一组条件。Few-NERD提供层级实体类型和少样本任务，可以用来考察
模型在已知粗类下区分细类、或在少量标注下学习新类别的能力；它并不使所有模型自动获得
新类型识别能力。其来源见本章文献说明中的\citetitle{ding2021fewnerd}。
评价时分别固定字符跨度、类型集合及嵌套规则；只比较边界、在给定边界上比较类型，
以及同时比较二者，是不同指标。类型正确但少取字符、边界完整但类型错误，都不计入
严格实体级真阳性。

\subsubsection{BIO--CRF与BiLSTM-CRF实体标注}
\label{sec:nlp-extraction-bio}

在平面实体规范下，可以把跨度转换为词元标签。BIO用\code{B-ORG}表示组织的首词元，
\code{I-ORG}表示其后续词元，\code{O}表示实体之外；BIOES再显式区分末尾和单词元实体。
设编码器给出$\bm{h}_i\in\mathbb{R}^{d}$，标签集合为$Y$，发射分数为
$\bm{u}_i=\bm{W}\bm{h}_i+\bm{b}\in\mathbb{R}^{|Y|}$。
独立词元分类只依据$\bm{u}_i$选择标签；条件随机场（conditional random field, CRF）
把相邻标签是否相容也纳入计算：
\begin{equation}
  s(\bm{y},x)=
  \sum_{i=1}^{n}u_{i,y_i}
  +\sum_{i=1}^{n}A_{y_{i-1},y_i},
  \qquad
  L_{\mathrm{seq}}=-s(\bm{y}^{*},x)
  +\log\sum_{\bm{y}\in Y_{\mathrm{legal}}}\exp s(\bm{y},x).
  \label{eq:nlp-extraction-crf}
\end{equation}
这里$\bm{y}^{*}$为标注序列，$\bm{A}$是含起始标签转移的得分矩阵；
结束转移可同样加入。非法转移掩码必须同时用于训练配分函数和解码：
\code{O}之后不能直接出现\code{I-ORG}，也不能从\code{B-LOC}直接接\code{I-ORG}。
通用CRF形式见第~\ref{sec:pgm-energy-models-crf-definition}节。

特征CRF的发射分数来自字形、词典和上下文特征；BiLSTM-CRF用双向长短期记忆网络
产生$\bm{h}_i$；预训练编码器标签头改变的也是这一得分来源。训练用标注标签序列
最小化式~\eqref{eq:nlp-extraction-crf}，反向更新编码器、分类头和转移参数。
推理复用上一章的合法转移与Viterbi过程，递推
$v_i(c)=u_{i,c}+\max_{c'}[v_{i-1}(c')+A_{c',c}]$，保存使最大值成立的先行标签，
在句尾回溯。时间复杂度为$O(n|Y|^2)$，并列得分按固定标签次序处理。

以教学切分“海棠楼／201／开放”为例，正确标签是
$[\texttt{B-LOC},\texttt{I-LOC},\texttt{O}]$，恢复“海棠楼201”。
$[\texttt{B-ORG},\texttt{I-ORG},\texttt{O}]$边界正确而类型错误；
$[\texttt{B-LOC},\texttt{O},\texttt{O}]$类型方向正确却漏了房间号。
若输入切成单一词元“海棠楼201开放”，普通词元标签甚至无法表示正确边界，
需要更细切分或字符级输出。分层标签序列可以覆盖某些嵌套规范，但必须给出各层实体的
对应和冲突处理规则，不能把一条合法BIO序列视为任意嵌套或不连续结构的完整解。

\subsubsection{跨度分类与双仿射起止评分}
\label{sec:nlp-extraction-biaffine}

当一个位置需要属于多个实体时，可以直接判断每个起止对。候选集合为
$C=\{(i,j):1\leq i\leq j\leq n\}$，每个候选有实体类型或“非实体”标签。
本书若为控制计算量设置$j-i+1\leq\ell$，这是额外的教学裁剪条件，
并非下述原方法必须使用的长度上限。未进入候选集合的金标准实体应计入候选漏检。

Yu等使用BERT、词与字符表示及BiLSTM形成上下文表示，再把同一位置分别投影为
起点和终点表示。其关键是对每个合法起止对作类别判断，而不是在NER输出上构造
依存树\scite{yu2020ner}
为写清维度，设$\bm{p}_i=f_{\mathrm{s}}(\bm{h}_i)$、
$\bm{q}_j=f_{\mathrm{e}}(\bm{h}_j)\in\mathbb{R}^{k}$，类型集合为$T$，
另加非实体类别$\varnothing$。类型$t$的双仿射分数为
\begin{equation}
  z_{i,j,t}=
  \bm{p}_i^\top\bm{U}_t\bm{q}_j+
  \bm{w}_t^\top[\bm{p}_i;\bm{q}_j]+b_t,
  \qquad
  p(t\mid i,j,x)=
  \frac{\exp z_{i,j,t}}{\sum_{t'\in T\cup\{\varnothing\}}\exp z_{i,j,t'}}.
  \label{eq:nlp-extraction-biaffine}
\end{equation}
其中$\bm{U}_t\in\mathbb{R}^{k\times k}$、$\bm{w}_t\in\mathbb{R}^{2k}$、
$b_t\in\mathbb{R}$；分号表示拼接。双线性项刻画两端的配对，
仿射项保留单独起点和终点的信息。普通跨度分类也可用
$[\bm{h}_i;\bm{h}_j;\operatorname{pool}(\bm{h}_{i:j});\bm{e}_{j-i+1}]$
输入多层感知机，$\bm{e}_{j-i+1}$是可学习的长度嵌入。

将金标准跨度置为对应类型，其余有效候选置为$\varnothing$，训练最小化
$L_{\mathrm{span}}=-\sum_{(i,j)\in C}\log p(t^{*}_{i,j}\mid i,j,x)$。
负候选通常远多于正候选；若另做负采样或类别重加权，须公布采样范围和权重，
并在完整候选集上评价。本式按每跨度单标签定义，同一跨度允许多类型的本体需要
改为多标签目标，不能悄悄用一个softmax丢弃标签。

\begin{example}[嵌套候选怎样保留]
\label{ex:nlp-extraction-nested}
教学词元序列为“海棠／大学／计算机／学院／发布／通知”，候选及其最高类别分数为
$(1,4,\mathrm{ORG},4.0)$、$(1,2,\mathrm{ORG},3.0)$、
$(2,5,\mathrm{ORG},2.0)$、$(3,4,\mathrm{ORG},1.5)$，
各候选的非实体分数均为0。为单独检查跨度恢复，此例只使用一个实体类别。
取$(1,4)$的$\bm{p}_1=(1,0)^\top$、$\bm{q}_4=(2,1)^\top$、
$\bm{U}_{\mathrm{ORG}}=\bm{I}_2$、$\bm{w}_{\mathrm{ORG}}=(0,0,1,0)^\top$、
$b_{\mathrm{ORG}}=0$，式~\eqref{eq:nlp-extraction-biaffine}得到$2+2=4$，
实体概率为$\exp(4)/(1+\exp(4))\approx0.9820$，该候选损失约为0.0181。

按最高类别分数降序扫描时，先保留$(1,4)$。嵌套规范保留其内部的$(1,2)$和$(3,4)$，
但排除交叉候选$(2,5)$，因为$1<2\leq4<5$；平面规范只保留$(1,4)$。
对应原文字符区间分别为$[0,9)$、$[0,4)$和$[4,9)$。
这是可复算的候选选择过程，不保证贪心结果使全部跨度的总分最大。
\end{example}

独立起止指针若只报告起点$\{1,3\}$和终点$\{2,4\}$，无法单凭集合判断
$(1,4)$是否应当存在。双仿射给每个起止组合独立评分，缓解了这种配对歧义，
但仍可能把错误组合判为实体。长度截断造成的漏检、表示造成的类别错误和解码造成的
冲突排除应分开统计。单个$(i,j)$也仍是连续跨度，不表达任意不连续提及。

\subsubsection{查询式阅读理解NER}
\label{sec:nlp-extraction-mrc-ner}

类型标签“组织”本身提供的信息很少。可以把标注指南改写成查询：
“找出学校、学院等组织名称”，并与通知一起送入编码器。这就是查询式阅读理解NER：
查询指定要找哪类对象，答案仍然必须是原文中的实体跨度。
组织查询可恢复“海棠大学”与“海棠大学计算机学院”，地点查询则应返回空集，
除非输入另外含有地点；不能因为大学有地理位置就把组织名称自动标成地点。

Li等的模型对每个词元分别预测“是否起点”“是否终点”，再判断两端能否匹配，
而不是只在整段文本上各选一个起点和终点\scite{li2020mrcner}
设类型查询$c_t$条件下的正文表示为$\bm{H}^{(t)}\in\mathbb{R}^{n\times d}$，
两个$d\times2$分类头产生逐行二分类概率。匹配头可写成
\begin{equation}
  p_{i,j}^{(t)}=\sigma\!\left(
  \bm{v}^{\top}[\bm{h}^{(t)}_i;\bm{h}^{(t)}_j]+b\right),
  \qquad
  L=\alpha L_{\mathrm{s}}+\beta L_{\mathrm{e}}+\gamma L_{\mathrm{match}},
  \label{eq:nlp-extraction-mrc}
\end{equation}
其中$\bm{v}\in\mathbb{R}^{2d}$，三个损失分别监督起点、终点和合法跨度配对，
权重由开发集选择。无此类型实体的文本提供全负边界与配对标签；多答案提供多个正位置。
训练更新共享编码器和各分类头，推理在预测的起止集合中枚举$i\leq j$，
经匹配头筛选并恢复字符偏移。它可以保留不同查询下的嵌套实体，同类型嵌套仍依赖匹配头。

固定类型查询、新类型文字描述、附加少样本示例提供的信息条件不同，比较时应分别记录。
描述中加入“学院”等例词能帮助限定标签含义，但不保证未见类别无需训练即可准确识别。
查询数增多也会增加重复编码成本。无答案、边界、类型和多答案配对仍需独立评价，
问题句式并没有替代这些约束。

\subsubsection{W2NER词元关系图}
\label{sec:nlp-extraction-w2ner}

跨度方法把中间全部文字视为实体的一部分。若标注允许不连续片段，就需要另外记录
哪些词元实际参与实体。W2NER用两类有向词元关系表达这一点：
NNW（next-neighboring-word）连接一个实体内部相邻的两个词元，
THW-$t$（tail-head-word）从该实体的尾词元指回首词元，并携带类型$t$。
这里的“相邻”指实体内部的次序相邻，不要求原文位置相邻；
来源见文献说明中的\citetitle{li2022w2ner}。

设正文有$n$个词元，模型输出$n\times n$关系网格，每格在
$\{\mathrm{NONE},\mathrm{NNW}\}\cup\{\mathrm{THW}\text{-}t:t\in T\}$中分类。
原方法由BERT和BiLSTM取得词元表示，经条件层归一化、距离与区域嵌入构造词元对表示，
用二维膨胀卷积交换相邻网格信息，再把双仿射分支和多层感知机分支的类别分数相加。
若两分支均输出$\bm{z}_{i,j}\in\mathbb{R}^{|T|+2}$，逐格softmax交叉熵即可
监督关系标签；NONE格也参与训练。网格占用随$n^2$增长，不能忽略长文档的候选成本。

教学句按词切成“左／和／右／肺／均／清晰”。本节明确采用省略共享中心词的标注规范，
把“左……肺”和“右肺”各记为一个身体部位提及。正标签为
$G_{1,4}=\mathrm{NNW}$、$G_{4,1}=\mathrm{THW}\text{-}\mathrm{BODY}$、
$G_{3,4}=\mathrm{NNW}$、$G_{4,3}=\mathrm{THW}\text{-}\mathrm{BODY}$，
其余格为NONE。解码从THW边的首位置出发，沿索引递增的NNW路径搜索到尾位置，
恢复路径$1\to4$和$3\to4$。前者输出字符片段$[0,1)$、$[3,4)$，
后者输出$[2,4)$；不能把“左和右肺”整体当成第一个实体。

缺失NNW边会断开片段，THW类型错误会改变整条路径的类别，多条路径共享边时还可能
组合出标注中不存在的实体。因此路径合法不等于实体真实；须核对每个恢复片段及其组合。
单词元实体可由对角线上的THW-$t$表示。不同语料是否把省略形式视作不连续实体必须
事先声明，不能通过改变标注规则提高结果。

\subsubsection{UIE结构化生成与模式提示}
\label{sec:nlp-extraction-uie}

当实体、关系和事件需要使用统一接口时，可以把抽取结构写成一个可解析序列。
UIE的结构化抽取语言（structured extraction language, SEL）用类型与片段的对应
表达定位，用下层角色与片段表达关联。例如教学结构
\code{((地点: 海棠楼201))}表示一个地点提及；增加下层关联可以表达关系或事件论元。
结构模式提示（structural schema instructor, SSI）把要识别的类型和关联名放在输入前缀，
说明本次任务要求抽取什么。原方法、预训练和模式接口见文献说明中的
\citetitle{lu2022uie}。

设模式提示为$c$，原文为$x$，金标准结构序列为$\bm{y}^{*}$。编码器处理$[c;x]$，
解码器在第$u$步由$\bm{g}_u\in\mathbb{R}^{d}$和
$\bm{W}_{\mathrm{v}}\in\mathbb{R}^{|V|\times d}$生成词表分布。
下游训练使用教师强制，即每步给出正确前缀：
\begin{equation}
  L_{\mathrm{gen}}=
  -\sum_{u=1}^{|\bm{y}^{*}|}
  \log p_\theta(y_u^{*}\mid \bm{y}_{<u}^{*},x,c).
  \label{eq:nlp-extraction-generation}
\end{equation}
UIE还通过文本与结构配对、纯结构和文本去噪数据进行预训练；
下游标注决定当前类型及片段监督，不能把预训练知识当作当前原文的事实来源。
推理用预测前缀逐步生成，遇结束符后解析括号、类型和片段，再对齐原文。

模式提示是条件，不是自动生效的硬约束。本书的教学抽取流程另外维护文法状态：
在类型位置只准输出类型词表，在片段位置可限制为原文可对齐的候选，括号必须闭合；
生成阶段过滤非法词元与生成后验证完整结构，是两个动作。某条路径若没有合法延伸，
返回失败或空结果，不通过补造实体凑齐结构。此处的文法和原文核验是教材补充要求，
不声称UIE已保证事实正确。

通知中“海棠楼201”若出现两次，复制字符串只能证明文字存在，不能唯一确定提及位置。
位置记录需要显式偏移或保留所有匹配候选，再借上下文消歧。输出一个原文没有的
“李老师”，即使括号和类型都合法，也应被判为无依据新增。
指令加示例可以替代显式SSI作为条件，验收仍须检查漏项、越界类型、定位歧义和原文支持。
生成模型的通用计算见第~\ref{chap:transformer}章。

\subsection{实体链接}
\label{sec:nlp-extraction-linking}

识别提及后，需要确定它是否对应已有的登记对象。\term{实体链接}（entity linking, EL）
把提及映射到给定版本知识库的身份标识；不能确认对应项时输出NIL。
为比较四种方法，构造两个同名“海棠大学”：$e_1$是教学大学，有“计算机学院”的描述；
$e_2$是教学小说中的虚构学校，有“故事人物”的描述。输入“海棠大学计算机学院发布通知”
应链接$e_1$，而“小说中的海棠大学迎来新角色”应链接$e_2$。

候选召回回答正确身份是否进入候选集合$C(m)$，重排回答在该集合中是否选对。
若正确身份已被遗漏，后续再准确的分类也不能恢复它。评价须声明知识库快照、
提及是否给定、候选规模和NIL比例；知识库之后新增了对象，并不使较早快照中的
NIL成为自动错误。

\subsubsection{别名字典与实体先验排序}
\label{sec:nlp-extraction-alias}

最直接的候选来源是登记名称、人工别名及历史链接记录。字典把“海棠大学”映射到
$\{e_1,e_2\}$；历史链接计数可估计$p(e\mid m)$。加入类型与上下文后，
一种明确的教学评分为
\begin{equation}
  s(m,e)=\log p(e\mid m)+
  \lambda\,\operatorname{sim}(x_m,d_e)+
  \eta\,\operatorname{type}(m,e),
  \label{eq:nlp-extraction-alias}
\end{equation}
其中$x_m$为提及上下文，$d_e$为实体描述，后两项分别为上下文相似度与类型兼容得分，
$\lambda,\eta$在开发集选择。字典本身不需要梯度训练；先验必须由训练期计数估计，
罕见别名可平滑，不能偷用测试链接频率。

若教学先验为$(0.8,0.2)$，小说语境的相似度为$(0,2)$，取$\lambda=1$且类型项相等，
两分数为$-0.2231$和$0.3906$，于是选$e_2$，而不是机械选择最常见的$e_1$。
在学院语境中把相似度换成$(2,0)$，则选$e_1$。若别名字典遗漏$e_2$，
这两段语境都无法在当前候选集内选到它。低分拒绝或前两名分差规则应另行校准；
字典查不到名字只能说明候选来源不足，不证明现实中没有该对象。

\subsubsection{BLINK双编码召回与交叉编码重排}
\label{sec:nlp-extraction-blink}

字面别名缺失时，还可根据上下文与知识库描述的语义相近程度检索。
BLINK分别编码带提及边界标记的上下文和“规范名称加描述”，再用联合编码重排候选。
其原论文的主要身份选择设定假定提及在库中有对应项，NIL须单独处理\scite{wu2020blink}
设双编码输出$\bm{u}_m,\bm{v}_e\in\mathbb{R}^{d}$，相似分数为
$s_{\mathrm{bi}}(m,e)=\bm{u}_m^\top\bm{v}_e$。
对于正确身份$e^+$及负候选集合$N_m$，训练目标为
\begin{equation}
  L_{\mathrm{bi}}=
  -\log\frac{\exp s_{\mathrm{bi}}(m,e^+)}
  {\sum_{e\in\{e^+\}\cup N_m}\exp s_{\mathrm{bi}}(m,e)}.
  \label{eq:nlp-extraction-blink}
\end{equation}
批内其他实体提供负例，检索得到的高分错误身份提供困难负例；同一真实身份不能误当负例。
交叉编码器将上下文与候选描述放在同一输入中，用
$s_{\mathrm{cross}}(m,e)=\bm{w}^{\top}\bm{g}_{m,e}+b$评分，
$\bm{g}_{m,e},\bm{w}\in\mathbb{R}^{d}$，再以正确候选的交叉熵训练重排头。

在学院语境的教学计算中，令$\bm{u}_m=(1,0.5)^\top$，
$\bm{v}_{e_1}=(0.8,0.4)^\top$、$\bm{v}_{e_2}=(0.9,0.5)^\top$，
点积分数为1.00与1.15。若只召回一项，正确身份$e_1$已经丢失；
召回两项后，假设交叉编码得分为$(2.4,0.2)$，则
$p(e_1\mid C(m),x)\approx0.9002$，恢复$e_1$。
这里的概率依赖候选集合，不等于库外拒绝概率。

知识库向量可预先计算，精确检索代价为每提及$O(|E|d)$，
近似近邻索引以召回损失换取查询成本下降；重排则要对$k$个候选分别联合编码。
实体描述或编码器变化后须更新索引。错误诊断应先检查召回覆盖，再检查重排，
另以库外样本训练或校准拒绝器，不能仅用候选softmax最大值宣称解决NIL。

\subsubsection{文档全局一致性重排}
\label{sec:nlp-extraction-global-linking}

同一文档的其他提及有时能帮助消歧。例如“海棠大学”和“计算机学院”的候选身份可能
分别有校园和小说语境的解释。局部分类各自最高分未必组成最合理的身份组合。
设文档提及为$m_1,\ldots,m_q$，候选赋值为$\bm{e}=(e_1,\ldots,e_q)$，
一种教学全局目标是
\begin{equation}
  \widehat{\bm{e}}=\argmax_{\bm{e}\in\prod_i C(m_i)}
  \left[\sum_i s(m_i,e_i)+
  \lambda\sum_{i<j}c(e_i,e_j)\right].
  \label{eq:nlp-extraction-global-linking}
\end{equation}
兼容分数$c$可来自给定知识库中的隶属关系、共现统计或标注文档学习的表示，
训练时必须说明哪一种来源；若联合学习，可用正确组合与错误组合的排序损失，
如$\max(0,1-S(\bm{e}^{*})+S(\bm{e}^{-}))$。

教学局部分数设为：大学的校园候选0.6、小说候选0.8，学院的校园候选0.9、
小说候选0.4。同领域兼容分数为0.7，跨领域为0，$\lambda=1$。
四种组合总分依次为校园--校园2.2、校园--小说1.0、小说--校园1.7、
小说--小说1.9；全局选择校园组合，改变了大学的局部首选。
小例可穷举，$q$个提及各有$k$候选时组合达到$k^q$，实际系统需要束搜索、
消息传递或整数优化等近似或受限求解，并说明搜索范围。

兼容并非越强越好。一篇讨论“现实大学与小说学校”的文章本来就可能跨领域，
强迫同领域会制造身份错误。全局重排也必须保留时间条件和NIL候选；
同一文档共同出现、同属一个共指簇，都不能脱离原文与库描述单独证明身份。

\subsubsection{生成式实体名称预测}
\label{sec:nlp-extraction-generative-linking}

知识库的身份还可以通过规范名称生成出来。GENRE把身份对应到唯一的文本名称，
以提及上下文为输入，训练自回归模型生成正确名称；它用知识库名称构造前缀树，
在解码时屏蔽不可能完成任何合法名称的延伸\scite{decao2021genre}
对于规范名称词元序列$\bm{y}_e$，评分为
$\log p_\theta(\bm{y}_e\mid x_m)=\sum_u\log p_\theta(y_{e,u}\mid\bm{y}_{e,<u},x_m)$，
训练使用式~\eqref{eq:nlp-extraction-generation}的序列负对数似然，
身份监督来自已链接提及。GENRE在屏蔽时保留原模型的合法词元对数概率，
不在剩余候选上重新归一化；束搜索仍是近似搜索，并不保证找到全局最高分名称。

教学知识库把两个对象规范化为“海棠大学（校园）”和“海棠大学（小说）”。
生成“海棠大学（”后，前缀树只允许通向这两个登记名称的词元路径；
若完成路径的概率分别为0.12和0.08，就选择前者并查表恢复$e_1$。
输入别名可以相同，输出名称必须能唯一映射身份。知识库若使用重名标题，
就需要消歧后缀或另存身份，不应把字符串当成天然唯一键。

新增名称要更新前缀树和映射表，但合法候选增加不意味着模型已经学会如何选择它。
无对应项时可另设NIL输出或拒绝判断，并单独评价。名称属于知识库只证明格式与候选域
合法，仍需核验它是否对应当前提及，不能用受约束生成替代身份判断。

\section{实体关系抽取}
\label{sec:nlp-extraction-relations}

找全组织和活动名称以后，还要判断谁举办了什么。\term{关系抽取}（relation extraction, RE）
把参与对象与关系表达绑定起来。例如，有向三元组
\[
  (\text{海棠大学},\text{主办},\text{读书会})
\]
表示前者主办后者，而不是反过来。
这里使用的组织是本节补充的教学材料，不由章首通知推断。二元关系有两个有序参与者，
多元事实还可能包含日期、地点和职责等角色；把三元组分开存储，不能丢掉它们属于同一
事实实例的依据。事件实例之间的先后或因果另见第~\ref{sec:nlp-extraction-event-relations}节。

句内或文档级说明允许读取的上下文，封闭模式或开放关系说明输出标签是否预先枚举，
二元或多元说明参与结构。这些是任务条件，不是互斥的模型类别。本节先确定对象怎样
进入同一个候选，再计算关系和证据，最后恢复带条件的记录。

\subsection{关系对象的绑定}
\label{sec:nlp-extraction-relation-binding}

给定文档内$M$个实体，最直接的二元候选有$M(M-1)$个有向实体对，不含自关系。
例如4个实体有12对；若本体只允许2个组织分别与1个活动构成“主办”关系，就只剩2对。
类型裁剪降低了计算量，也把上游类型错误变成不可恢复的候选漏检。句距限制同样会遗漏
跨句事实，应同时报告候选召回率和候选内分类结果。

对象可以是某次提及，也可以是文档共指簇或库中身份，三者不能混用。
“海棠大学于4月18日在海棠楼201主办读书会”至少需要组织、活动、时间、地点。
如果另一句写“晨光书院于4月19日在松柏楼举办研讨会”，从全文独立选最高分的
组织、日期和地点，就可能拼出原文没有的记录。多元关系应保留共同实例键及各角色
到该实例的绑定；若采用二元关系重述，也要通过同一个实例节点连接这些字段。

实体给定时，分类器只解决候选关系判断；联合抽取还必须定位参与者。
以下两种方法使用同一补充教学句，词元依次为
“海棠大学／和／晨光书院／共同／主办／读书会”。
金标准为$(1,\text{主办},6)$和$(3,\text{主办},6)$，编号在此简写单词元提及；
它们共享客体，但有不同主体。原文字符跨度分别为$[0,4)$、$[5,9)$和$[13,16)$。

\subsubsection{CasRel级联主体--关系--客体抽取}
\label{sec:nlp-extraction-casrel}

CasRel先找主体，再在指定主体条件下，按每个关系分别寻找客体。
这样，关系不只是两个已经选好的名字之间的分类标签，而是从主体到客体集合的映射。
其方法出处为\citeauthor{wei2020casrel}（\citeyear{wei2020casrel}），完整信息见本章文献说明。
设编码器输出$\bm{h}_i\in\mathbb{R}^{d}$。主体起止头为
$p_i^{\mathrm{sh}}=\sigma(\bm{w}_{\mathrm{sh}}^\top\bm{h}_i+b_{\mathrm{sh}})$与
$p_i^{\mathrm{st}}=\sigma(\bm{w}_{\mathrm{st}}^\top\bm{h}_i+b_{\mathrm{st}})$，
每个位置都是二分类，不要求一句话只有一个主体。对于主体跨度$(a,b)$，取
$\bm{v}_s=(\bm{h}_a+\bm{h}_b)/2\in\mathbb{R}^{d}$，关系$r$的客体头为
\begin{equation}
  p_{i,r}^{\mathrm{oh}\mid s}
  =\sigma\!\left(\bm{w}_{r,\mathrm{oh}}^\top
       (\bm{h}_i+\bm{v}_s)+b_{r,\mathrm{oh}}\right),
  \qquad
  p_{i,r}^{\mathrm{ot}\mid s}
  =\sigma\!\left(\bm{w}_{r,\mathrm{ot}}^\top
       (\bm{h}_i+\bm{v}_s)+b_{r,\mathrm{ot}}\right).
  \label{eq:nlp-extraction-casrel}
\end{equation}
所有$\bm{w}$均为$d$维；关系条件体现在不同的客体头参数中。
主体标签由全部金标准主体构造，客体标签则由“当前主体及当前关系”的全部客体构造。
令$\operatorname{BCE}(y,p)=-y\log p-(1-y)\log(1-p)$，
训练求主体两类边界损失与各已标注主体条件下客体两类边界损失之和，
更新共享编码器和标记头。若每步只抽样一个金标准主体，须说明抽样方案；
不能把该主体以外的真实客体不加条件地混入当前标签。

在教学句中，主体起点、终点正位置均为$\{1,3\}$。给定$s=(1,1)$，
“主办”客体两端的正位置均为6；给定$s=(3,3)$也如此。
假设两次客体起止概率为$(0.9,0.8)$与$(0.85,0.9)$，教学阈值取0.5，
便恢复两条三元组。第一对正边界的损失为$-\log0.9-\log0.8\approx0.3285$；
这只是一对边界的贡献，完整损失还包含负位置及主体预测。

解码先枚举主体起点，为其匹配最近的合法主体终点，再按主体逐关系恢复客体。
这类最近终点规则计算简单，却不适合表示同一起点同时对应多个嵌套终点的全部情形。
主体漏检会使其全部关系失去后续预测机会；正确主体下的客体漏检是另一个错误来源。
输出仍须检查方向、类型和字符位置，两个主体共享客体也不能自动说明它们互相合作。

\subsubsection{TPLinker词元对联合标记}
\label{sec:nlp-extraction-tplinker}

TPLinker（\citeyear{wang2020tplinker}，见本章文献说明）把绑定依据直接放到词元对上：
实体的首尾是否相接，某关系的主体首词元与客体首词元是否对应，
主体尾词元与客体尾词元是否对应。
对$i\leq j$，以
$\bm{g}_{i,j}=\tanh(\bm{W}[\bm{h}_i;\bm{h}_j]+\bm{b})\in\mathbb{R}^{k}$
构造词元对表示，其中$\bm{W}\in\mathbb{R}^{k\times2d}$。
上三角的$n(n+1)/2$个位置按固定顺序铺成handshaking序列；
它是共享的索引布局，不是先后执行的关系推断链。

EH--ET头标记实体首尾，输出“无连接／有连接”。
每个关系另有SH--OH和ST--OT两个头，分别标记主体与客体的首部、尾部连接。
因为只存上三角，后两类头各含“无连接、正向、反向”三个标签；
反向标签记录主体位置大于客体位置的情形，不能在排序索引时丢掉方向。
因此有$2|R|+1$个标记头，各头将$\bm{g}_{i,j}$线性投影到自己的标签空间，
对全部有效格作交叉熵训练。

六词元教学句有21个格。实体头的正格为$(1,1)$、$(3,3)$、$(6,6)$；
“主办”的首部连接和尾部连接均在$(1,6)$、$(3,6)$取正向标签。
推理先由实体头恢复跨度集合$E$。对每个关系$r$，枚举
$s=(a,b),o=(c,d)\in E$，仅当SH--OH给出有向$a\to c$，
且ST--OT给出同方向的$b\to d$时，才恢复$(s,r,o)$。
本例得到两个共享客体的三元组。若只有$(1,6)$的首部连接，却没有对应尾部连接，
就不能从高分首部边单独输出关系。

CasRel依赖先选出的主体，TPLinker依赖三类连接的共同一致性。
后者避免了推理时逐主体调用客体头，却增加了二次规模的词元对输出；
独立连接还可能错误组合，尤其在重复边界和密集重叠中。实体边界正确、关系连接正确、
完整三元组正确应分别检查，不能把所有连接边自由拼接为任意多元事实。

\subsection{关系与证据的预测}
\label{sec:nlp-extraction-relation-evidence}

文档级关系往往需要多个句子。构造三个教学句：
$S_1$“海棠大学设有机器学习中心”；
$S_2$“该中心负责读书会的筹备”；
$S_3$“海棠楼201配有投影设备”。
若任务允许“所属机构承担下属单位的筹备职责”这一明确声明的推理规则，
可以由$S_1,S_2$提出“海棠大学参与筹备读书会”的候选；若本体只标直接陈述关系，
就不应自动补出这条关系。文档背景与可用推理规则共同决定监督目标。

DocRED将文档实体、关系与证据句组织为可评价对象，并区分人工标注与远程监督数据；
来源见文献说明中的\citetitle{yao2019docred}。
远程监督通过已有知识库与文本对象对齐产生候选标签，实体对在知识库中有某关系，
不意味着当前文档表达了它。未标注关系也未必是真负例，须说明标注覆盖及评价口径。
相关来源噪声可通过本章第~\ref{sec:nlp-extraction-merge}节的来源模型理解，
但来源可靠性不能弥补原文根本未表达当前事实的问题。

以下计算以给定文档实体为条件。证据必须支持具体对象、方向和条件：
“不负责筹备”“如获批准则负责筹备”“去年负责筹备”不能恢复为同一条无条件现时关系。
UIE式输出可复用第~\ref{sec:nlp-extraction-uie}节的接口，
但需增加主客体、关系及证据字段，不能把共享生成器当成共享监督含义。

\subsubsection{句内实体对分类与依存路径特征}
\label{sec:nlp-extraction-sentence-re}

对于已绑定的有向提及对$(m_s,m_o)$，可以先拼接词汇和位置特征$\bm{\phi}$，
以$\bm{z}=\bm{W}\bm{\phi}+\bm{b}$计算关系分数。
特征包含两个实体的词、类型、相对距离以及最短依存路径上的词和弧向。
例如补充句“海棠大学主办读书会”中的路径为
“海棠大学$\leftarrow$主办$\rightarrow$读书会”；
主语和宾语的方向信息帮助区分谁主办谁，不能仅保留一个无向词袋。

编码器方法在两个提及前后加入不同边界标记，得到
$\bm{h}_s,\bm{h}_o,\bm{h}_{\mathrm{sent}}\in\mathbb{R}^{d}$，
再令$\bm{\phi}=[\bm{h}_s;\bm{h}_o;\bm{h}_{\mathrm{sent}}]$，
$\bm{W}\in\mathbb{R}^{(|R|+1)\times3d}$，额外一类表示NA。
单标签规范用softmax交叉熵训练；同一对允许多个关系时，应改为逐关系二分类或
下一节的多标签阈值目标。实体顺序交换会改变输入，不应默认关系对称。

教学得分按“主办、支持、NA”取$(2,0,-1)$，
则$p(\text{主办})=\exp2/(\exp2+1+\exp(-1))\approx0.8438$，
正确标签的负对数似然约为0.1698，解码恢复有向“主办”关系。
此数值仅说明分类计算；证据仍需回到句子及提及位置。
若句子变成“海棠大学并不主办读书会”，最短路径可能仍包含相同的核心三词，
路径外的否定修饰必须进入表示。句法错边、对象绑定错误和关系分类错误要分开诊断，
句内路径也不能覆盖全部跨句事实。

\subsubsection{ATLOP文档实体对预测}
\label{sec:nlp-extraction-atlop}

一个实体可能在多处被称呼，不同实体对又需要不同上下文。
ATLOP分别用多提及聚合、局部上下文池化和自适应阈值处理这些问题
\scite{zhou2021atlop}
设文档编码为$\bm{H}\in\mathbb{R}^{n\times d}$，每次实体提及前插入标记。
实体$e$有提及标记表示$\bm{h}_{e,1},\ldots,\bm{h}_{e,M_e}$，
逐维对数和指数聚合为
\begin{equation}
  [\bm{h}_e]_u=\log\sum_{j=1}^{M_e}\exp[\bm{h}_{e,j}]_u,
  \qquad u=1,\ldots,d.
  \label{eq:nlp-extraction-atlop-entity}
\end{equation}
这一表示可以汇集不同提及的突出维度，却不是对提及身份作新的证明；
把两个同名对象误并为一个实体，会直接污染聚合结果。

设最后一层注意力有$H_{\!a}$个头，先在每个头内平均同一实体全部提及标记的注意力行，
得到$\bm{A}_s,\bm{A}_o\in\mathbb{R}^{H_{\!a}\times n}$。
局部上下文池化对两个实体共同关注的位置赋权：
\begin{equation}
  q_i=\sum_{h=1}^{H_{\!a}}A_{s,h,i}A_{o,h,i},
  \qquad a_i=\frac{q_i}{\sum_{j=1}^{n}q_j},
  \qquad \bm{c}_{s,o}=\bm{H}^{\top}\bm{a}\in\mathbb{R}^{d}.
  \label{eq:nlp-extraction-atlop-context}
\end{equation}
这里先在同一个头内相乘，再对头求和；先把各头平均后相乘会产生不同头之间的交叉项，
不是相同的计算。掩去填充位置后归一化，数值实现还需防止分母下溢。
该权重表示模型对实体对的关注，不是人工证据概率。

融合后令$\bm{z}_s=\tanh(\bm{W}_s\bm{h}_s+\bm{W}_{cs}\bm{c}_{s,o})$，
$\bm{z}_o$类似，四个投影矩阵均为$d\times d$。
将两向量各分为$K$组，每组$d/K$维，关系$r$的分数为
\begin{equation}
  \ell_r=\sum_{g=1}^{K}
  \bm{z}_{s,g}^{\top}\bm{U}_{r,g}\bm{z}_{o,g}+b_r,
  \qquad
  \bm{U}_{r,g}\in\mathbb{R}^{(d/K)\times(d/K)}.
  \label{eq:nlp-extraction-atlop-score}
\end{equation}
分组降低双线性参数量。参数区分主体和客体，因此相同局部上下文不使关系自动对称。
除实际关系外，模型还输出一个阈值类别TH，其分数$\ell_{\mathrm{TH}}$也依赖当前实体对。

令$P$为该对的正关系集合，$N=R\setminus P$为训练中按所选监督口径处理的负关系。
ATLOP的目标把正关系推到TH以上、负关系推到TH以下：
\begin{equation}
  \begin{aligned}
    L_{+}&=-\sum_{r\in P}
      \log\frac{\exp\ell_r}
      {\exp\ell_{\mathrm{TH}}+\sum_{r'\in P}\exp\ell_{r'}},\\
    L_{-}&=-\log\frac{\exp\ell_{\mathrm{TH}}}
      {\exp\ell_{\mathrm{TH}}+\sum_{r'\in N}\exp\ell_{r'}},
    \qquad L_{\mathrm{AT}}=L_{+}+L_{-}.
  \end{aligned}
  \label{eq:nlp-extraction-atlop-loss}
\end{equation}
特别要注意，$L_+$的每个分母包含全部正类和TH，不是分别对每个正类与TH做二分类。
$P$为空时$L_+=0$，$L_-$仍训练TH。推理取
$\widehat{R}_{s,o}=\{r\in R:\ell_r>\ell_{\mathrm{TH}}\}$，
空集表示NA；它不是只取softmax最高类，也不是固定sigmoid阈值0.5。

\begin{example}[从局部上下文到多关系输出]
\label{ex:nlp-extraction-atlop}
为单独检查计算，使用四个位置、一个注意力头和二维表示。
令$\bm{H}$的四行依次为$(1,0)$、$(0,1)$、$(2,1)$和$(1,2)$，
前两行是主体的两个提及标记表示，后两行是客体的提及标记表示。
主体和客体分别逐维聚合，得到
\[
  \begin{aligned}
    \bm{h}_s&=(\log(1+\mathrm{e}),\log(1+\mathrm{e})),\\
    \bm{h}_o&=(\log(\mathrm{e}+\mathrm{e}^{2}),
                \log(\mathrm{e}+\mathrm{e}^{2})).
  \end{aligned}
\]
两实体平均后的注意力行取
\[
  \bm{A}_s=(0.4,0.3,0.2,0.1),\qquad
  \bm{A}_o=(0.1,0.2,0.3,0.4).
\]
逐位置乘积为$(0.04,0.06,0.06,0.04)$，其元素之和为0.2。
归一化后，权重向量为$\bm{a}=(0.2,0.3,0.3,0.2)$。
再用同一$\bm{H}$得到上下文向量$\bm{c}_{s,o}=(1.0,1.0)$。

用固定参数使融合向量均为$(0.5,0.5)$：取实体投影为零，
上下文投影为$\operatorname{arctanh}(0.5)\bm{I}_2$。
取$K=1$、偏置为零，并令$r_1,r_2,r_3,\mathrm{TH}$的双线性矩阵依次为
$4\bm{I}_2,2\bm{I}_2,-2\bm{I}_2,\bm{0}$，
式~\eqref{eq:nlp-extraction-atlop-score}给出分数$(2,1,-1,0)$。
若$P=\{r_1,r_2\}$，则
$L_+=2\log(\exp2+\exp1+1)-3\approx1.8152$，
$L_-=\log(1+\exp(-1))\approx0.3133$，总损失约2.1285。
解码同时保留$r_1,r_2$，排除$r_3$。这里的简化参数只用来贯通计算，
不表示训练模型忽略实体表示；实际关系名称、方向和证据必须另按原文核验。
\end{example}

局部上下文解决“这个实体对该看哪里”，阈值类别解决“这个实体对该保留哪些标签”，
两者不互相替代。注意力集中的词可能只是背景，多关系输出也可能同时犯错；
实体与提及给定条件下的成绩不能代替从原文定位、共指、绑定开始的端到端结果。

\subsubsection{EIDER证据增强与推理融合}
\label{sec:nlp-extraction-eider}

ATLOP没有直接训练“哪些句子足以支持关系”。EIDER在共享文档编码器上增加
实体对条件的证据预测，并在推理时结合完整文档与证据伪文档的关系分数。
其方法出处为\citeauthor{xie2022eider}（\citeyear{xie2022eider}），完整信息见本章文献说明。
设文档有$J$句，第$j$句的表示$\bm{s}_j\in\mathbb{R}^{d}$
由句内词元向量逐维对数和指数池化获得。给定实体对上下文$\bm{c}_{s,o}$，
论文的证据概率可写为
\begin{equation}
  p_j^{s,o}=\sigma(\bm{s}_j^\top\bm{W}_v\bm{c}_{s,o}+b_v),
  \qquad \bm{W}_v\in\mathbb{R}^{d\times d}.
  \label{eq:nlp-extraction-eider-evidence}
\end{equation}
它按实体对选句，不是为每个关系类别分别选一套句子。
EIDER论文在此先平均注意力头，再对两个实体的注意力逐位置相乘并归一化；
公开实现则采用与式~\eqref{eq:nlp-extraction-atlop-context}等价的逐头乘积聚合，
证据头还有分组与多头实现。复现时应选定一个口径，不混合声称为同一公式。

设$D^+$为至少具有一个正关系的训练实体对集合，$y_j^{s,o}$为该对第$j$句的证据标签。
使用标准二元交叉熵，
\begin{equation}
  \begin{aligned}
    L_{\mathrm{evi}}
      &=-\sum_{(s,o)\in D^+}\sum_{j=1}^{J}
      \left[y_j^{s,o}\log p_j^{s,o}
       +(1-y_j^{s,o})\log(1-p_j^{s,o})\right],\\
    L&=L_{\mathrm{RE}}+\lambda L_{\mathrm{evi}}.
  \end{aligned}
  \label{eq:nlp-extraction-eider-loss}
\end{equation}
实际可对有效项取平均，但要保持归约和$\lambda$的定义一致。
不能把全部NA实体对的全部句子都当作负证据加入原方法的证据目标。
关系头使用自适应阈值损失，两个任务共同更新编码器。
论文联合目标取$\lambda=1$，公开代码使用0.1；
上式BCE按损失定义及作者代码书写，原文相关公式正项漏排的对数不应沿用。

证据可以来自人工标注，也可以来自启发式规则。EIDER的规则按先后后备方式寻找
同句共现、共指相关句及共同桥接实体形成的句组，并非对所有来源无区别求并集。
桥接实体是分别与两个端点有关联的第三实体，多个候选时按规则选取；
它仍可能带来不支持目标关系的句子。使用规则版本时，还应区分规则直接选证据的
推理流程与学习型证据头预测的流程，不能把启发式监督称为独立人工核验。

推理对每个实体对按阈值选择证据句，保持原句序，拼成该对的伪文档。
需要同时重建提及位置及“伪文档位置到原文位置”的映射。
两路关系模型分别产生完整文档分数$\ell_r^O$与证据文档分数$\ell_r^E$，
各自减去各自的TH分数：
\begin{equation}
  S_r^O=\ell_r^O-\ell_{\mathrm{TH}}^O,\qquad
  S_r^E=\ell_r^E-\ell_{\mathrm{TH}}^E,\qquad
  p_r^{\mathrm{fuse}}=\sigma(S_r^O+S_r^E-\tau).
  \label{eq:nlp-extraction-eider-fusion}
\end{equation}
$\tau$是在开发集拟合的单一标量，原融合不是两路概率的凸加权，
也不是把两路隐藏向量拼接后重新训练分类器。
以$p_r^{\mathrm{fuse}}>0.5$判正等价于$S_r^O+S_r^E>\tau$。
空证据或证据文档缺少实体时，本书实现应返回完整文档结果并标记回退；
这是明确补充的保护规则，不是原论文对证据完整性的保证。

对上述实体对，设三句的证据标签为$(1,1,0)$。
预测概率取$(0.8,0.7,0.2)$时，证据损失之和为
$-\log0.8-\log0.7-\log0.8\approx0.8030$。
取0.5选句阈值，伪文档包含$S_1,S_2$。假设某个允许的筹备关系两路原始
“关系、TH”分数分别为$(0.6,0.2)$和$(0.8,0.1)$，
则$S^O=0.4,S^E=0.7$。取教学融合阈值$\tau=0.3$，
得到$\sigma(0.8)\approx0.6900$并输出该关系。
若证据遗漏使第二路变成$S^E=-0.5$，融合降为$\sigma(-0.4)\approx0.4013$。
可见高分选句和正确关系并不相互保证：应分别检查证据遗漏、背景误选和关系误判，
尤其要复核否定、方向及所允许的跨句推理规则。

\subsubsection{ReVerb开放关系抽取}
\label{sec:nlp-extraction-reverb}

封闭模式必须预先定义“主办”等标签；开放抽取则先保留文本中的关系短语。
ReVerb依靠英语词性模式及词汇约束识别关系，再找左右论元，最后给抽取置信度
\scite{fader2011reverb}
其句法候选可概括为$V$、$VP$或$VW^{*}P$：
$V$为动词，$P$为介词、动词小品词等规定位置，$W$为方法允许的中间词性集合。
按规则保留最长匹配并处理相邻匹配；多词关系还要检查其词汇模式是否在语料中
连接过足够多的不同论元对，以排除把具体论元误吞进关系的过度特定短语。
这一步的参数与词汇表来自语料统计，不是下游封闭关系标签监督。

为展示原英语机制，明确把补充通知译写为教学句
“Haitang University hosts the reading group.”
关系候选为“hosts”，左侧最近合规则名词短语是“Haitang University”，
右侧是“the reading group”，得到原文表达三元组。
论元规则会排除某些关系代词、疑问词及存在句的“there”等不合格左论元；
不能任意取左右两个专名。方法另用标注的正确、错误抽取训练逻辑回归置信度，
例如$p(\text{有效}\mid\bm{\phi})=\sigma(\bm{w}^{\top}\bm{\phi})$，
特征刻画关系与论元的句法和抽取条件，推理按阈值保留候选。

封闭模式可将“hosts”归一为“主办”，开放结果则保留“hosts”及其位置。
“does not host”中的否定、时间状语和条件从句若被丢弃，短三元组就可能改变原意，
需要下一节的限定字段承接。ReVerb的上述资源与规则面向英语；
中文“主办”示例只是任务对应，不能声称直接运行英语词性模式就已完成中文抽取。
开放标签扩大了表达范围，却不保证论元边界、同义归一或限定条件自动正确。

\subsection{关系记录的恢复与核验}
\label{sec:nlp-extraction-relation-records}

模型分数还不是可用记录。恢复结果至少包含对象标识或未链接提及、关系、限定条件、
文档版本和证据位置。教学记录可以表示为
$\langle s,r,o,\text{polarity},\text{time},\text{condition},\text{source}\rangle$；
否定关系的极性为负，不宜把它删除后留下同样的正关系。
以JSON Schema检查对象字段的类型、必填项、枚举值和数组结构，
只说明记录符合接口\scite{wright2022jsonschema}
日期字符串还需真正解析：若使用\code{format: date-time}，
须启用格式断言或另做校验，因为某些实现默认只把格式当注解。

复制片段应通过原始文档版本和字符区间恢复。重复字符串返回多个候选位置，
必须结合主体、客体及上下文消歧；轻微改写可以用编辑距离产生候选，
却不能把最小距离匹配当成已经确认的来源。
开放关系经别名表映射到规范关系时，同时保存原措辞和方向：
“由海棠大学主办”应调换表面论元次序后恢复组织到活动的主办边。
类型兼容检查可以拒绝“日期主办组织”，但不能判断一条类型兼容的边是否真实。

结构生成阶段的文法约束阻止非法延伸，生成后的JSON Schema检查完整对象，
SHACL则对图中的类型、属性、范围和基数进行约束验证
\scite{knublauch2017shacl}
例如活动允许多个主办方，就不能为了方便给该关系强设最大基数1。
这几种机制检查不同层次的结构，不是三种事实证明。
证据--记录支持判定可把原文作为前提、规范记录的文字表达作为假设，
使用蕴含或句对模型筛查，再独立检查对象身份、数值、否定、时间和条件。
“未参与主办”与“参与主办”只有很小字符串差异，却不语义等价；
“负责举办”与“主办”可能在当前本体中等价，须由预先声明的映射决定。

精确匹配按固定对象键、规范关系和条件字段比较；
语义等价评价另声明允许的别名、时间表达和方向改写，保留逐项判定记录。
流水线的候选遗漏、联合模型的连接误组、生成模型的无依据补项，最终都应落回
相同的对象、关系、条件、来源检查。格式通过不能抵消内容错误，
而给定正确实体条件下的关系分数也不能掩盖真实输入中的身份错绑。

\section{事件抽取}
\label{sec:nlp-extraction-events}

“海棠大学主办读书会”给出关系，“读书会从4月17日改至4月18日”则描述了一次改变。
\term{事件}（event）是任务本体要求记录的一次发生、计划或状态变化；
\term{事件提及}（event mention）是文中对它的一次表达。
同一事件可以被多句重复提及，同样的事件类型也可以在不同日期多次发生。
本节依次形成事件提及、角色论元和共指簇，保留各层产出。

语义角色标注描述谓词及其参与者，事件本体却可能只关心“会议改期”及其旧时间、新时间、
活动等角色。两者的目标和角色集合不必相同。不能把句法主语一概填为组织者，
也不能把语义角色分析得到的全部谓词都自动当成当前任务应抽取的事件。

\subsection{事件提及与类型识别}
\label{sec:nlp-extraction-event-detection}

\term{触发词}（trigger）是用于定位事件提及的词或短语。
在本节教学本体中，“会议延期”的“延期”触发改期事件；
“拟延期”中的触发仍可定位为“延期”，但状态是计划，不是已确认发生。
事实性可用发生、否定、假设、计划等属性表达，具体标签须以任务标注为准。
没有事实性监督的基准，不能仅凭事件类别得分声称已完成状态判断。

MAVEN提供大规模、多类型的事件检测任务，适合考察类型覆盖和不均衡，
但它是数据与任务设置，不是一种检测算法；来源见文献说明中的
\citetitle{wang2020maven}。
评价需区分给定触发跨度的类型分类与从原文开始的定位加分类。
对于“本周活动安排如下”这类没有典型触发词的表达，若本体仍要求抽取活动，
须规定标题、句子或记录段落作为锚点；只输出类型名称无法区分文中多个实例。

\subsubsection{触发词BIO标注与事件跨度分类}
\label{sec:nlp-extraction-trigger-tagging}

对平面触发结构，可把实体BIO标签换成事件类型标签，沿用
式~\eqref{eq:nlp-extraction-crf}的发射、转移和负对数似然；
解码还原触发跨度及类型。跨度方法则枚举$(i,j)$，
使用式~\eqref{eq:nlp-extraction-biaffine}或端点拼接分类器，
在事件类型与非触发类别之间判断。训练时未标触发的候选充当负例，
这要求明确语料是否充分标注目标类型，不能把本体外事件当成“现实中没有事件”。

教学句“会议拟延期”按“会议／拟／延期”切分时，第三词元标为
\code{B-RESCHEDULE}，恢复字符$[3,5)$；“拟”的字符位置$[2,3)$提供计划状态证据。
若另有状态标签，可用触发表示及周围上下文$\bm{g}\in\mathbb{R}^{d}$
预测$\bm{p}_{\mathrm{status}}=\operatorname{softmax}(\bm{W}_{f}\bm{g}+\bm{b}_f)$，
并在有标注样本上加入状态交叉熵。将“会议延期”改为“会议拟延期”，
触发类型仍可判为改期，字符位置却从$[2,4)$移到$[3,5)$，状态也须改为计划。

同一跨度若在本体中允许多个事件类型，单一BIO类别或跨度softmax不能完整表达。
可为每个类型设置独立sigmoid并以多热标签训练，或使用明确的多层结构；
恢复时保存多个“跨度、类型”记录，再核验它们是多个事件实例还是同一实例的多标签。
不应为了适应模型接口静默删除一个标签。触发检测完成后仍缺参与者和时间等论元，
本节产出不能直接称为完整事件抽取。

\subsubsection{DMCNN动态多池化触发分类}
\label{sec:nlp-extraction-dmcnn}

同一句话可能包含几个事件，用整个句子的一个最大值向量难以判断哪个特征属于当前候选。
DMCNN（\citeyear{chen2015dmcnn}，见本章文献说明）在卷积之后，
按候选位置划分池化区域，使相对位置进入特征。
设每个词的词向量与相对候选位置嵌入拼为$\bm{x}_i\in\mathbb{R}^{d_x}$，
宽$w$的卷积窗口经$\bm{W}_c\in\mathbb{R}^{F\times wd_x}$产生
$\bm{c}_i=f(\bm{W}_c[\bm{x}_{i:i+w-1}]+\bm{b}_c)\in\mathbb{R}^{F}$。
触发检测只有一个候选锚点$t$，因此沿它分成两段：
\begin{equation}
  [\bm{p}_{L}(t)]_u=\max_{i\leq t}[\bm{c}_i]_u,\qquad
  [\bm{p}_{R}(t)]_u=\max_{i>t}[\bm{c}_i]_u.
  \label{eq:nlp-extraction-dmcnn}
\end{equation}
窗口对齐和边界填充需统一。池化只考虑有效对齐位置；
当$t=n$使右区域为空时，本节教学实现约定$\bm{p}_R(t)=\bm{0}\in\mathbb{R}^{F}$，
不对空集直接取最大值，也不让填充位置参与最大值比较。
拼接两个$F$维池化向量与候选附近的局部词表示，
经线性层、softmax预测事件类型或NONE，以候选触发标注的交叉熵训练。
原方法在论元分类时有触发和论元两个锚点，才划成三段，
不能把那一规则照搬到本节的单锚点检测。

教学句为“会议／延期／后／报名／暂停／一天”。
为只看池化作用，设某一卷积通道在六个对齐位置的响应为$(1,4,2,3,5,0)$。
候选“延期”在$t=2$时，池化为$(4,5)$；“暂停”在$t=5$时为$(5,0)$。
全句最大池化都会得到5，动态池化保留了这个强特征位于候选哪侧的信息。
真实计算还会因相对位置嵌入改变卷积响应，本例固定响应是为了隔离池化差异。

该模型仍需枚举候选并训练大量NONE样本。候选切分不含完整触发短语时，
池化不能恢复被排除的跨度；局部最高响应来自另一事件时也可能混淆类型。
因此位置条件帮助分辨上下文，并不保证多事件绑定，更不自动补全论元。

\subsubsection{查询式事件检测}
\label{sec:nlp-extraction-event-query}

事件类型定义可以作为检测条件，而不只作为输出端的一个整数。
本节用教学实现说明：查询“哪些词表达了活动时间的更改”，
输入查询与章首通知，预测“改至”的起止位置；
换成“哪些词表达了活动取消”，则应输出空集。查询描述提供了标签含义，
比较时必须记录是否额外给了定义、例子和领域词表。

设查询$q_t$和原文经联合编码得到$\bm{H}^{(t)}\in\mathbb{R}^{n\times d}$，
按第~\ref{sec:nlp-extraction-mrc-ner}节的逐位置边界二分类与配对头评分。
金标准提供“类型查询、触发跨度列表”，无该类型的查询全标负；
训练最小化边界及匹配损失，推理枚举超过阈值的起止对并恢复原文锚点。
这里特意采用多答案的逐位置二分类，不与下文EEQA跨位置softmax混为同一目标。

若同一通知包含两次更改，查询应恢复两个有位置区别的触发；
若允许无显式触发事件，需把候选单位改为本体规定的句子或段落锚点，
并用对应锚点标注训练。查询形式本身不能创造锚点，也不证明所述变化已实际发生。
类型、位置和有监督的状态属性仍分别验收，新类型文字可作为条件，
但不保证仅凭名称就能得到可靠的零样本检测。

\subsubsection{Text2Event序列到结构生成}
\label{sec:nlp-extraction-text2event}

当提及和论元需要联合产生时，可以直接生成事件树。Text2Event把事件类型、触发和
角色论元组织为树，再按深度优先次序线性化；目标结构和课程学习安排见文献说明中的
\citetitle{lu2021text2event}。
事件节点关联触发片段，其子节点关联角色和论元片段，同级信息按规定的原文次序排列。
训练以原文为输入、线性化结构为目标，使用
式~\eqref{eq:nlp-extraction-generation}的教师强制负对数似然。

在本节教学改期模式下，章首通知可转为如下可读结构，
括号中的名字是本体定义，不是原模型特定词元编码：
\begin{quote}
（改期 改至（活动 机器学习读书会）（原时间 4月17日）\\
（新时间 4月18日（周五）14:00））
\end{quote}
实际序列应使用独立结构标记与原文字符的转义规则，避免把日期里的普通括号当树边界。
这里“改至”定位改期提及，新时间与原时间属于同一次改期，不能分别附着到相邻的
另一场活动。没有目标事件时也应有明确的空结构目标，例如空括号。

Text2Event的课程先学习触发、独立论元等较简单子结构，再训练完整事件树；
这样改变训练难度次序，不改变最终必须恢复的完整目标。
其模式约束解码根据当前树位置，用前缀树限制事件类型、角色名及原文片段等候选，
同时控制括号结构。解码器每步对词表打分，屏蔽非法延伸后做束搜索，
结束后解析树并对齐片段。没有合法延伸时应报告失败，不生成未经允许的角色补齐树。

UIE着重统一结构语言与模式提示接口，Text2Event的这里则是面向事件树的序列化、
课程和约束解码。两者都可能生成合法但绑错对象的结构：
原文有两个相同日期时，日期字符串本身不能唯一确定证据位置。
本节分别检查触发和类型，角色、位置及完整记录继续接受下一节的核验，
不能因输出已经长成一棵树而省略这些指标。

\subsection{事件论元填充}
\label{sec:nlp-extraction-arguments}

确定一次事件后，才能问它的时间、地点和参与者是什么。
\term{事件论元}（event argument）是以特定角色参与该事件的对象或表达，
角色集合由本体定义，不由某个句法位置直接决定。
为检查绑定，补充构造一份通知：
“4月18日14:00，机器学习读书会在海棠楼201举行。4月19日10:00，
数据讨论会在松柏楼102举行。两场活动均由海棠大学组织。”
第二场活动和组织信息仅用于本节教学，不改写章首共同通知。

WikiEvents用于考察文档级论元抽取及信息更充分的论元表达。
“它”可以是直接提及，“机器学习读书会”是其共指链中更明确的表达；
选哪个作为答案取决于标注规范，不能在评价时随意替换。
先确认当前任务是否给定触发词，再说明候选覆盖、多角色、多论元和论元缺席的处理。
以下方法都应保存事件键；两个事件允许共享“海棠大学”，
却不能因此共享日期或被合并为一场活动。

\subsubsection{事件--论元跨度分类}
\label{sec:nlp-extraction-argument-spans}

给定事件$e$，从实体提及或一般文本跨度形成候选集合$C_e$。
后者可覆盖“4月18日14:00”等尚未被NER列为实体的表达。
设事件表示$\bm{t}_e$和论元表示$\bm{a}_{i,j}$均为$d$维，
事件类型、句距等附加嵌入为$\bm{u}_{e,i,j}\in\mathbb{R}^{k}$。
一种教学分类头为
\begin{equation}
  \bm{z}_{e,i,j}
  =\bm{W}[\bm{t}_e;\bm{a}_{i,j};
             \bm{t}_e\odot\bm{a}_{i,j};\bm{u}_{e,i,j}]+\bm{b},
  \quad \bm{W}\in\mathbb{R}^{(|A|+1)\times(3d+k)}.
  \label{eq:nlp-extraction-argument-spans}
\end{equation}
$A$是角色集合，额外一类为无角色；$\odot$表示逐元素乘积。
单角色规范对每个事件--候选对作softmax交叉熵，多角色规范改用独立二分类。
类型模式掩码可屏蔽该事件不允许的角色，但掩码必须符合真实本体。

训练标签按事件键构造。上述通知中“海棠楼201”对读书会是地点正例，
对数据讨论会却是无角色；“海棠大学”对两者都可以是组织者。
若地点与无角色两类得分对前者为$(2,0)$、后者为$(-1,1)$，
地点概率分别为$\sigma(2)\approx0.8808$与$\sigma(-2)\approx0.1192$。
这说明候选文本相同，事件条件不同会得到不同角色判断。
推理选择角色并恢复“事件、角色、字符跨度”，如本体要求单一开始时刻，
还需在该事件内部处理多个冲突候选，而非逐个高分全部写入。

跨句裁剪可能排除末句的组织者。以给定候选为条件的角色正确率不显示这一漏检，
必须另报告候选覆盖和完整事件指标。共享跨度图可以联合优化事件、实体和角色边，
模式提示也可进入编码，但都不免除事件绑定。
语义角色标注提供的谓词结构可以作特征，仍需按事件角色重新标注、训练与验收。

\subsubsection{EEQA角色问答式抽取}
\label{sec:nlp-extraction-eeqa}

EEQA（\citeyear{du2020eeqa}，见本章文献说明）把事件论元抽取转为角色条件的阅读理解，
问题可包含事件类型、触发及角色，从原文直接找答案，无需前置实体候选。
例如针对第一场“举行”询问“机器学习读书会何时举行”，
对第二场则明确询问数据讨论会。只问“何时举行”而不区分触发上下文，
会使两场活动的答案都看似合理。

设联合输入的表示为$\bm{H}\in\mathbb{R}^{n\times d}$，
包含一个专用CLS位置。起止向量$\bm{w}_s,\bm{w}_e\in\mathbb{R}^{d}$
对全部有效输入位置评分，再分别跨位置归一化：
\begin{equation}
  P_s(i)=\frac{\exp(\bm{h}_i^\top\bm{w}_s)}
               {\sum_{u=1}^{n}\exp(\bm{h}_u^\top\bm{w}_s)},
  \qquad
  P_e(j)=\frac{\exp(\bm{h}_j^\top\bm{w}_e)}
               {\sum_{u=1}^{n}\exp(\bm{h}_u^\top\bm{w}_e)}.
  \label{eq:nlp-extraction-eeqa}
\end{equation}
有答案训练以正确边界最小化$-\log P_s(i^*)-\log P_e(j^*)$，
同一问题的多个金标准论元提供多个边界监督项；没有该角色时起止标签均为CLS。
对CLS的负对数似然会提高无答案概率，不能误写成无答案时要降低CLS概率。
它与MRC-NER逐位置二分类不同：两个边界头各自的全部位置概率之和为1。

推理不能只取一个最高分跨度，否则会遗漏多个同角色论元。
EEQA枚举正文中的合法$i\leq j$及允许长度，排除
$P_s(i)<P_s(\mathrm{CLS})$或$P_e(j)<P_e(\mathrm{CLS})$的候选，
再计算
\begin{equation}
  \delta(i,j)=P_s(\mathrm{CLS})+P_e(\mathrm{CLS})-P_s(i)-P_e(j).
  \label{eq:nlp-extraction-eeqa-threshold}
\end{equation}
在开发集上扫描阈值，保留$\delta(i,j)$不大于该阈值的候选，
因此一次问题可恢复多个论元。此阈值按开发集任务指标选择，
不是ATLOP为每个实体对预测的TH类别。

为手算，假设有效位置仅有“CLS／4月18日／14:00／海棠楼201”，
起点概率为$(0.10,0.55,0.05,0.30)$，
终点概率为$(0.10,0.10,0.65,0.15)$，且最大答案长度为2。
日期时间跨度$(2,3)$的$\delta=0.2-0.55-0.65=-1.00$；
地点跨度$(4,4)$为$0.2-0.30-0.15=-0.25$。
取教学阈值$-0.60$只保留前者，恢复“4月18日14:00”。
它的边界训练损失为$-\log0.55-\log0.65\approx1.0286$。
完整输入还会包含其他位置，本例概率是为说明解码构造的小分布。

针对地点、参与者分别发问时必须重新计算条件分布，不能重复使用上述时间问题概率。
直接位置预测能找出NER漏掉的时间表达，但也可能从第二场活动取到高分答案。
无答案判断只回答“是否输出”，不保证已输出论元属于正确事件。

\subsubsection{事件模板条件生成}
\label{sec:nlp-extraction-event-template}

WikiEvents论文中的条件生成方法把角色槽位模板与文档一起输入模型，
围绕给定触发生成填槽后的模板\scite{li2021wikievents}
模板先依据事件类型写好，如本节教学模板
“\code{<arg>}于\code{<arg>}在\code{<arg>}举行”。
这些槽位依固定模板位置分别映射活动、时间、地点，
不能仅因占位符写法相同就认作同一角色。输入还标记当前事件触发，
区分文档里两个“举行”。

对读书会，目标文本为“机器学习读书会于4月18日14:00在海棠楼201举行”；
对讨论会则为“数据讨论会于4月19日10:00在松柏楼102举行”。
训练使用目标模板序列的条件负对数似然。某角色缺失时保留\code{<arg>}，
一个槽位的多个论元按约定连接词组织，随后按模板语法解析回角色列表。
推理把可输出词元限制到输入词元集合，输入包含模板及文档；
这一限制减少库外词元，却不能保证生成内容是原文中的一个连续跨度，
也不能保证各词元来自同一次事件。

如果末句用“该校”指代海棠大学，信息更充分的答案可以选择其已核实共指表达，
但要保留原始论元提及、共指连接及所选表述的来源。
两个角色槽位都生成“海棠大学”时，仍须核验是否确指同一对象；
两个日期字符串相同也不证明事件相同。
WikiEvents是评价基准，模板条件生成是具体方法，不能用基准名称替代计算过程。
给定触发的论元结果、端到端完整事件结果与信息更充分表达的评价应分别报告。

\subsection{事件共指与合并}
\label{sec:nlp-extraction-event-coreference}

一篇通知和它的转发可能描述同一次读书会，连续两周的读书会却是不同实例。
\term{事件共指}（event coreference）把表达同一次事件的提及归为一簇。
候选阻塞可以使用事件类型、触发词义、时间范围或可靠活动编号，
减少全体事件提及两两比较；时间缺失时不能直接把候选拒绝，否则会漏掉省略日期的转述。

设事件对特征$\bm{\phi}_{a,b}$包含触发语境、已绑定参与者、时间与地点兼容性、
事实性及来源，模型预测
$p_{a,b}=\sigma(\bm{w}^{\top}\bm{\phi}_{a,b})$，
用已标共指对的二元交叉熵训练。编码器或双线性头可以替换线性评分，
但正例必须指同一次事件，不能仅是相同类型。
推理按分数从高到低处理候选，受约束凝聚聚类每次合并两个簇前检查全部簇级限制，
例如互斥的可靠活动编号。不同日期是否硬冲突取决于是否允许同一活动改期，
应先辨认旧时间、新时间与不同场次。

构造三条记录：$a$为编号\code{ML-0418}的通知，
$b$为省略编号但明确转引该通知的报道，
$c$为另一个编号\code{ML-0425}的下周读书会。
教学相似分数取$p_{a,b}=0.92,p_{b,c}=0.86,p_{a,c}=0.10$，
阈值0.8。简单对高分边取传递闭包会把三者合为一簇；
受约束聚类先合并$a,b$，再因可靠编号冲突拒绝与$c$合并，
保存$\{a,b\}$和$\{c\}$。若编号本身未核实，则应转为待核验，
而非把相似分数当成编号真伪的裁决。

事件先行项排序也可为每个提及在此前事件与“无先行项”中选择，
共享表示的联合模型则同时训练抽取和共指；二者改变误差传播与优化方式，
并不取消簇内一致性检查。可靠记录键可直接归并，但必须声明键在来源和时间范围内唯一。
保留各次提及与其来源后，才能回查错误合并。实体共指连接参与者，
事件共指连接发生实例，知识库身份连接登记对象，三者分别验收。

\section{事件关系抽取}
\label{sec:nlp-extraction-event-relations}

两次事件并未合并为一次，不意味着它们彼此无关。报名先于活动、改期由场地冲突引起、
讨论是会议的一个环节，分别表达时间、因果和组成联系。
MAVEN-ERE在共同事件输入上组织共指、时间、因果及子事件关系，
有助于比较这些任务的相互影响；来源见文献说明中的\citetitle{wang2022mavenere}。
EventRelBench则通过句级、文档级等条件下的选择题评估事件关系理解
\scite{gong2025eventrelbench}
给定事件对并选择答案，不等于从原文定位全部事件、再恢复所有关系。
本节固定事件身份后讲关系计算，端点定位和共指错误仍应在端到端评价中计入。

\subsection{时间表达与时间关系}
\label{sec:nlp-extraction-temporal}

时间信息有三层对象：“明天”是待归一化的表达，某个日期属于哪次事件是关联，
两事件是否先后或重叠是关系。把“明天”转成日期，不代表已经找到了正确事件；
把两次事件各自的开始时间找对，也未必知道它们的结束时间和区间关系。
章首通知给出了活动开始时刻，却没有时长，不能自行补出结束时间。

本节保留明确时点、有限区间及尚未确定的参照条件。
日期粒度的“4月18日”表示一个日历日范围，不必等同于当天零点发生；
“4月16日18:00前”是严格早于该时刻的截止条件，
不能写成包含等号的闭区间，也不能压缩为“4月16日前”。

\subsubsection{HeidelTime与SUTime时间归一化}
\label{sec:nlp-extraction-time-normalization}

规则时间系统先识别日期、时刻、时长或周期表达，再结合语言资源和参照点归一化。
HeidelTime将模式、归一化资源与文本类型条件用于时间抽取，
SUTime也以规则和组合机制识别、解析时间表达\scite{strotgen2010heideltime}
SUTime的原始介绍以英语资源为主，不能把工具名等同于完整的中文覆盖
\scite{chang2012sutime}
这里的中文规则是明确的教学适配：用数字与日期词表匹配“$k$天后”，
从上下文选出参照日期$d_0$，计算$d_0+k$个日历日。

假设补充句发布于2025年4月15日：“明天18:00前报名，三天后14:00开始。”
“明天”匹配相对日规则并以文档日期为参照，得到
\code{2025-04-16T18:00:00+08:00}。
“三天后”则从发布日期增加三个日历日，得到
\code{2025-04-18T14:00:00+08:00}。
与章首通知核对，4月18日确为周五，报名条件是时间$t$严格满足$t<$截止时刻。
如果另一条发表于4月17日的咨询写“明天下午到”，其“明天”指4月18日，
但“下午”不足以证明会在14:00前到达，还需检查校外来宾是否已经报名。

参照不总是文档日期。“读书会4月18日举行，三天后公布纪要”的后半句以活动日为参照，
得到4月21日；机械加到发布日期会错成4月18日。
规则应输出表达位置、值、粒度、参照来源和未决字段，
再由事件关联步骤把这些时间绑定到报名、活动或纪要发布。
缺少参照时可保留$\text{date}=d_{\mathrm{ref}}+3\text{天}$，
不为了生成标准格式猜一个日期。日历加法也要正确处理月末、闰年和时区。

\subsubsection{事件对时间关系分类}
\label{sec:nlp-extraction-temporal-classifier}

给定事件$e_i,e_j$，模型输入包含两事件的触发或锚点、时间表达及允许上下文，
形成$\bm{g}_{i,j}\in\mathbb{R}^{d}$。
关系头$\bm{W}\in\mathbb{R}^{|T|\times d}$给出
$p(t\mid e_i,e_j,x)=\operatorname{softmax}(\bm{W}\bm{g}_{i,j}+\bm{b})_t$，
以标注的有向时间标签交叉熵训练。标签集$T$可以是先于、后于、同时、包含、被包含、
重叠和信息不足，也可以使用更细的区间代数，必须随语料规范声明。

对“报名截止时点、活动开始时点”这一有向对，已知时间为4月16日18:00和4月18日14:00，
两者相差44小时，因此输出前者先于后者。若教学“先于、后于、信息不足”得分为
$(2,-1,0)$，先于概率约0.8438；交换事件顺序时应输出逆关系。
这不意味着“所有报名行为”都有同一个时间，也不表示整个活动持续44小时。
时间关联对象必须先固定，不能在分类后再改写端点含义。

“三天后公布纪要”提供明确的参照关系，即使没有绝对发布日期也可判活动先于纪要发布；
“随后另有安排”若没有可定位事件，则连端点是否存在都尚待解决。
时间信息不足应保留可能关系或专门标签，未标注边不能统一当作无关系。
局部分数可能产生$a$先于$b$、$b$先于$c$、$c$又先于$a$的矛盾，
因此还需要检查全局结构。

\subsubsection{Allen区间约束与全局关系恢复}
\label{sec:nlp-extraction-allen}

Allen用13种基本关系描述两个非退化时间区间的位置关系：
先于、相接、重叠、同始而先结束、严格位于内部、同终而晚开始，
这六种及其逆关系，再加相等\scite{allen1983intervals}
设区间$I=[a,b)$与$J=[c,d)$，其中$a<b,c<d$。
例如先于要求$b<c$，相接要求$b=c$，重叠要求$a<c<b<d$，
严格位于内部要求$c<a<b<d$。同一个事件只有开始时刻时，
不能直接套入完整区间关系；应保留未知终点或使用另外声明的点时间规则。

一条不确定关系可表示为基本关系集合$R_{i,j}$。
关系组合$R_{i,j}\circ R_{j,k}$给出两条边允许的$i,k$关系集合，
路径一致性反复收缩
\begin{equation}
  R_{i,k}\leftarrow R_{i,k}\cap
      (R_{i,j}\circ R_{j,k}).
  \label{eq:nlp-extraction-allen-propagation}
\end{equation}
例如$\mathrm{before}\circ\mathrm{before}=\{\mathrm{before}\}$，
$\mathrm{meets}\circ\mathrm{meets}=\{\mathrm{before}\}$，
$\mathrm{during}\circ\mathrm{during}=\{\mathrm{during}\}$。
若$a$先于$b$、$b$先于$c$，而$R_{a,c}=\{\mathrm{before},\mathrm{overlaps}\}$，
更新后只剩before；若原集合只有after，交集为空，暴露冲突。
这一步使用已声明语义规则，不产生新的文本证据。

要兼顾模型分数，可以为基本关系设置二元变量$x_{i,j}^{r}$，目标为
\[
  \max_{\bm{x}}\sum_{i\ne j}\sum_r s_{i,j}^{r}x_{i,j}^{r}.
\]
要求每对选一个关系、逆关系一致：
$\sum_r x_{i,j}^{r}=1$，$x_{i,j}^{r}=x_{j,i}^{r^{-1}}$。
三元组合约束可写成
\begin{equation}
  x_{i,j}^{r}+x_{j,k}^{s}
  -\sum_{t\in r\circ s}x_{i,k}^{t}\leq1.
  \label{eq:nlp-extraction-temporal-ilp}
\end{equation}
当前两条边均被选中时，第三条必须落入组合允许集合。
若只处理严格先后，还可用次序变量或环约束保证无环。
未知关系不应为了满足“每对一个标签”被任意写成唯一事实：
可只求可行补全并保留原未知状态，或在声明的粗标签体系中显式保留不确定输出。

教学三事件的候选先后边权为$a\to b:9$、$b\to c:8$、$c\to a:6$，
反方向得分均为0，要求在三者确有先后顺序的条件下求一致赋值。
局部取最大产生环；枚举六个顺序，$a<b<c$总分17最高，
它舍去权重6的边，并由约束得到$a<c$。若事件实际上可重叠，
就必须扩大关系集合，不能沿用这个先后限定算例。
对任意Allen关系网络，局部或三元路径一致性并非充分可满足性条件；
完整一致性还可能需要端点约束求解，也不必得到唯一总时间线。

\subsection{因果关系}
\label{sec:nlp-extraction-causal}

“天气变差后会议取消”表达先后，“因天气变差会议取消”表达文本中的因果主张。
因果抽取先辨认这样的主张及其方向，不是在观测数据上估计干预效应。
即使原文明确说“因”，也只能据此记录作者作出的因果陈述，
不能自动证明外部世界的因果机制。

显式候选可由“因为、导致、由于”等连接词、因果谓词及依存路径产生，
再将原因和结果绑定到已识别的事件实例。
设两事件表示为$\bm{u},\bm{v}\in\mathbb{R}^{d}$，
关系$r$的一个教学双仿射头为
$z_r=\bm{u}^{\top}\bm{W}_r\bm{v}
+\bm{w}_r^{\top}[\bm{u};\bm{v}]+b_r$，
其中$\bm{W}_r\in\mathbb{R}^{d\times d}$。
输出标签可为原因指向结果、结果指向原因、明确无因果及信息不足；
以有向事件对标注作交叉熵训练。SVM使用词汇或路径特征，
编码器使用上下文表示，改变的是分数来源，不是因果标签的含义。

教学标注把“因天气变差会议取消”记为天气变化$\to$取消，
证据为完整因果从句；“天气变差后会议取消”若无更多材料则记为信息不足，
不能仅因先后就作为因果正例或明确否定例。
“天气变差使交通中断，也导致会议取消”中，交通中断与会议取消可能共享原因；
不能据此无条件补出交通中断导致会议取消。
未标注对也不统一视为负例，应按语料覆盖筛选训练和评价样本。

时间、因果、组成可共享事件编码器并优化
$L=\lambda_tL_t+\lambda_cL_c+\lambda_sL_s$，但各自保留标签、方向和监督掩码。
“原因应不晚于结果”只能在所用因果语义及时间粒度允许时作为约束，
不能无条件应用于持续过程或预期原因。
证据条件问答可以询问“文中给出的取消原因是什么”，
从答案跨度恢复原因事件并核对结果绑定；受约束生成则输出
“原因事件键、结果事件键、主张状态、证据”，过滤不存在的端点。
这些接口仍需逐项验证，外部知识只在任务明确允许时使用，
并把文本主张与外部支持分别保存。

\subsection{组成与子事件关系}
\label{sec:nlp-extraction-subevents}

补充教学安排“本次会议由专题报告和分组讨论两个环节组成”，
支持报告、讨论是这次会议的子事件。
它不同于事件共指：报告不是整场会议；也不同于类型上下位：
“学术会议是一类会议”谈的是类别，不是实例组成。
一个午休活动即使在会议的时间范围内，也未必是会议的一部分，时间包含不是充分条件。

给定父候选$p$与子候选$c$，把两者锚点、参与者及共同活动上下文编码为$\bm{g}_{p,c}$，
用线性或双仿射头预测“包含子事件、属于父事件、无组成关系”，
以标注的有向实例对交叉熵训练。不确定或未覆盖的边应单独掩码或保留状态。
推理按类型兼容和同一活动依据选择候选，恢复父$\to$子边及来源；
“会议”这个词在两处出现，不能在尚未消歧时把两个实例的所有环节合并。

若所选定义要求严格组成关系无环，可对选边变量$z_{p,c}$作整数优化：
最大化$\sum_{p,c}w_{p,c}z_{p,c}$，以拓扑次序变量$u_p$约束
\begin{equation}
  u_p+1\leq u_c+M(1-z_{p,c}),\qquad
  0\leq u_p\leq N-1,\quad M=N,
  \label{eq:nlp-extraction-subevent-ilp}
\end{equation}
其中$N$是事件节点数；选中边时强制父在子之前，这里的次序只是图拓扑次序，
不是事件时间先后。实际求解还要加类型和已核实实例背景约束。
允许一个环节参与多个上层活动时，结构应为有向无环图，而非强制单父树。

教学候选为会议$\to$报告得分3.0、报告$\to$讨论得分2.0、
讨论$\to$会议得分0.2。局部阈值0.1会选出环，
无环最大权子图舍弃0.2的边，保留总分5.0。
但原教学安排把报告和讨论列为并列环节，因此“报告$\to$讨论”仍缺语义支持；
结构合法只清除了环，不能替代证据核验，最终应依原文保留会议到两个环节的边。
传递闭包可记录已确认直接边蕴含的间接组成，传递约简可删去冗余展示边，
两者都不能创造原文没有的新依据。评价时必须说明目标是直接边还是包含传递边的关系集。

\section{知识记录的构建}
\label{sec:nlp-extraction-records}

抽取系统可能从每份通知都正确读出字段，汇总后仍然出错：把两场活动误并，
把一条转发当成新的独立证据，或用旧通知覆盖更正日期。
知识构建因此不只是把三元组装入图中，还要确定哪些记录可比较、
哪些值可以合并、冲突为何产生，以及新材料改变了哪一条已有主张。
知识图谱是表达这些连接的一种形式，带条件和来源的表格记录也可以承担同样的核验责任。

图~\ref{fig:nlp-extraction-records}把本章各层产出连成处理流程。
其中旧通知是额外教学构造，现通知仍是章首2025年4月15日的材料；
保留两者不等于认为活动将在两个日期各举行一次。
版本化的是同一活动的排期主张，改期关系提供了新旧值之间的解释。
\begin{figure*}[htbp]
  \centering
  % Chapter-local TikZ diagram; final width about 168 mm, no external assets.
% Solid arrows carry records. The dashed return arrow requests evidence checks.
\begin{tikzpicture}[
  x=1mm,y=1mm,
  font=\figurefont\fontsize{9}{11}\selectfont,
  box/.style={rectangle,align=left,inner sep=2mm,line width=0.85pt,
    draw=BookRule,fill=white,text width=42mm,minimum height=19mm},
  input/.style={box,draw=FigureInput,fill=FigureInput!8},
  model/.style={box,draw=FigureModel,fill=FigureModel!8},
  output/.style={box,draw=FigureOutput,fill=FigureOutput!8},
  flow/.style={line width=1pt,draw=BookInk,-{Latex[length=2mm]}},
  feedback/.style={line width=0.85pt,draw=FigureModel,dashed,
    -{Latex[length=2mm]}},
  heading/.style={font=\figurefont\bfseries\fontsize{9}{11}\selectfont,
    text=BookInk,anchor=west}
]
  \node[heading] at (0,17) {A\quad 每条事实保留可回查的依据};
  \node[input] (text) at (23,0) {
    \textbf{原始文本}\\
    文档ID、版本、发布日期\\
    保留未经覆盖的原文
  };
  \node[input] (mention) at (84,0) {
    \textbf{提及定位}\\
    字符跨度、类型\\
    词元到原文的映射
  };
  \node[model] (binding) at (145,0) {
    \textbf{对象绑定}\\
    共指簇、实体身份\\
    活动与角色归属
  };
  \node[model] (candidate) at (145,-34) {
    \textbf{关系或事件候选}\\
    对象、角色、极性\\
    时间、条件、证据句
  };
  \node[model] (check) at (84,-34) {
    \textbf{记录核验}\\
    模式、身份、原文支持\\
    限定条件与来源
  };
  \node[output] (record) at (23,-34) {
    \textbf{带来源记录}\\
    规范键、字段与值\\
    有效时间、写入时间
  };
  \draw[flow] (text.east) -- (mention.west);
  \draw[flow] (mention.east) -- (binding.west);
  \draw[flow] (binding.south) -- (candidate.north);
  \draw[flow] (candidate.west) -- (check.east);
  \draw[flow] (check.west) -- (record.east);
  \draw[feedback] (check.north) -- node[right,text=FigureModel] {回查} (mention.south);
  \draw[BookRule,line width=0.7pt] (0,-52) -- (168,-52);
  \node[heading] at (0,-60) {B\quad 同一活动改期：保留旧值，不把两次主张当成两场活动};
  \node[input,text width=75mm,minimum height=21mm] (old) at (40,-80) {
    \textbf{旧通知：4月14日发布}\\
    活动编号 ML-0418\\
    原排期为4月17日，未给时刻
  };
  \node[input,text width=75mm,minimum height=21mm] (new) at (127,-80) {
    \textbf{更正通知：4月15日发布}\\
    同一活动改至4月18日14:00\\
    “原定”与“现改至”连接新旧值
  };
  \node[output,text width=163mm,minimum height=25mm] (versions) at (84,-119) {
    \textbf{同一活动键，两个排期版本（年份均为2025年）}\\
    旧值：4月17日；排期主张有效期 $[\text{4月14日},\text{4月15日})$；保留旧来源\\
    新值：4月18日14:00；排期主张自4月15日起有效；关联更正来源\\
    系统若4月16日才读到更正，写入日为4月16日，有效起日仍为4月15日
  };
  \draw[BookInk,line width=1pt] (old.south) -- (40,-98) -- (84,-98);
  \draw[BookInk,line width=1pt] (new.south) -- (127,-98) -- (84,-98);
  \draw[flow] (84,-98) -- (versions.north);
\end{tikzpicture}

  \caption{从文本证据到知识记录。上部保存定位、身份、关系或事件以及来源字段，
  核验失败后返回相关定位或绑定步骤。下部两条通知均为教学构造；
  旧通知发布于4月14日，新通知发布于4月15日，本例约定更正从发布日生效。
  排期主张的有效期与活动开始时刻分开记录，系统写入日期也不决定事实的新旧。}
  \label{fig:nlp-extraction-records}
\end{figure*}

\FloatBarrier
\subsection{记录对齐}
\label{sec:nlp-extraction-alignment}

对齐先解决“字段说的是否是同一种东西”，再解决“记录说的是否是同一个对象”。
一份表用“活动编号、举行时间、地点”，另一份用“event\_id、start、room”，
可以通过模式映射对应；但“报名截止”不能因为同为日期就映射到start。
单位、时区、粒度和关系方向也要规范化，并保留原值与转换规则。
“14:00”与“下午2点”在同一日、同一时区可对应，
“4月18日”与“4月18日14:00”则是不同精度的信息，不能无条件视为完全相同。

规范化后的记录按可靠编号、日期块或名称词项形成候选。
阻塞避免两个来源的全部$m\times n$对比较，却可能漏掉别名、日期缺失和改期记录，
因此候选覆盖必须单独评价。模式映射成功不意味着身份匹配成功；
一个汇总记录对应多次活动时，可以有一对多或多对多关系，应保留对应类型和不确定性。
纯文本表格可直接抽取字段；页面图像还需文字识别与版面恢复，
那些前置错误不能由记录匹配模型默默修补。

\subsubsection{Fellegi--Sunter记录匹配}
\label{sec:nlp-extraction-fellegi-sunter}

Fellegi--Sunter把两个记录是否同指视作匹配$M$与非匹配$U$之间的判别，
以比较向量在两种条件下出现的概率形成证据\scite{fellegi1969linkage}
比较向量$\bm{\gamma}$可包含名称编辑距离的分档、名称词项相似度、地点和日期一致性。
核心似然比为
\begin{equation}
  \Lambda(\bm{\gamma})=
  \frac{P(\bm{\gamma}\mid M)}{P(\bm{\gamma}\mid U)}.
  \qquad
  \log\Lambda(\bm{\gamma})
    =\sum_{f\in O}
       \log\frac{m_f(\gamma_f)}{u_f(\gamma_f)}
       \quad\text{（字段条件独立时）}.
  \label{eq:nlp-extraction-fellegi}
\end{equation}
$O$为可比较字段，$m_f,u_f$可由已核验匹配、非匹配记录对估计，
未标数据也可在明确的混合模型下迭代估计。分项相加依赖给定匹配状态的字段条件独立，
不是原始似然比定义的必然性质；名称与简称高度相关时，重复计分会夸大证据。

教学参数取名称相似的$m/u=0.9/0.1=9$，地点一致的$m/u=0.8/0.2=4$。
比较“机器学习读书会”和已经列入别名表的“ML读书会”，若地点相同而日期缺失，
则$\Lambda=9\times4=36$，$\log\Lambda\approx3.5835$。
缺日期时，在假定缺失机制不额外提供身份信息的条件下对该字段边际化，
贡献为1，而不是把它当成不一致。
若只剩名称证据，似然比降为9。

为展示三段判定，规定教学阈值$\Lambda<3$拒绝、$\Lambda>20$接受，
其间交待核验，则36接受、9待核验。
如果可比较且不涉及更正的活动日期确实冲突，取日期不一致似然比
$0.05/0.95=1/19$，总比变为$36/19\approx1.8947$而拒绝。
阈值须按误匹配、漏匹配成本和目标错误率校准，上述数值不是通用阈值。
似然比也不是后验概率；给定先验匹配率$\pi$，
后验为$\pi\Lambda/(\pi\Lambda+1-\pi)$。

日期更正是这一例子的边界：同一活动改变日期，不构成身份不同的直接证明。
应先根据活动键、原定与现改等字段判明比较的是什么主张，
再使用合适的比较模型。缺失概率若与来源或匹配状态有关，
还需建模缺失模式，不能沿用“缺字段贡献1”的简化条件。

\subsubsection{监督匹配与受约束分配}
\label{sec:nlp-extraction-supervised-matching}

监督分类器可以学习字段交互，例如名称相近但学校不同通常不是同一活动。
输入规范字段、缺失掩码及相似特征$\bm{\phi}_{i,j}$，输出
$p_{i,j}=\sigma(f_\theta(\bm{\phi}_{i,j}))$，
以已核验同指对和非同指对的二元交叉熵训练。
负例既包括随机不同记录，也包括同名、同场所等困难负例；
训练和测试应按真实对象或时间分割，避免同一活动的不同转发跨集合泄漏。
字段双编码可先召回名称、描述语义相近的候选，再联合编码重排，
其候选遗漏与候选内错误仍分开检查。

如果两个名单确实都“一条记录对应一次活动”，且任务要求一对一对应，
可以在候选边效用$s_{i,j}$上求
\begin{equation}
  \max_{\bm{x}}\sum_{i,j}s_{i,j}x_{i,j},
  \qquad
  \sum_jx_{i,j}\leq1,\quad
  \sum_ix_{i,j}\leq1,\quad
  x_{i,j}\in\{0,1\}.
  \label{eq:nlp-extraction-assignment}
\end{equation}
效用可由校准后的匹配对数几率减去接受成本构造，
未召回的边禁止选择，允许未匹配时添加相应虚拟节点。
不能只对每一行独立取最高分，因为多行可能占用同一个目标。

教学效用矩阵为
$\bm{S}=\left(\begin{smallmatrix}0.90&0.80\\0.85&0.20\end{smallmatrix}\right)$，
两个来源各有两条记录，并在此例中要求完整一对一分配。
局部首选都指向第一列；对角分配总分1.10，交叉分配总分1.65。
用匈牙利算法求解时，把效用换成成本$\bm{C}=\bm{1}-\bm{S}$，
行、列减去最小值不改变完整分配的最优解：
\begin{equation}
  \begin{pmatrix}0.10&0.20\\0.15&0.80\end{pmatrix}
  \ \longrightarrow\
  \begin{pmatrix}0&0.10\\0&0.65\end{pmatrix}
  \ \longrightarrow\
  \begin{pmatrix}0&0\\0&0.55\end{pmatrix}.
  \label{eq:nlp-extraction-hungarian}
\end{equation}
现在可选不同行列的零位$(1,2),(2,1)$，回到原成本得到$0.20+0.15=0.35$，
即效用1.65的交叉分配。
一般情况下，若零位尚不能形成完整匹配，就以最少行列覆盖全部零位，
从未覆盖位置减去最小未覆盖值，在双重覆盖位置加回该值，
再沿交替路径增广匹配，直到得到足够多的独立零位。

这个优化只在声明的一对一条件下合理。若一条记录是全年汇总、另一来源按每期分行，
强制一对一会丢掉真实对应；若候选分数都低，也不应为了填满矩阵强配。
全局最优是“给定候选、效用和约束下”的最优，不是身份事实的保证。

\subsection{记录合并与冲突处理}
\label{sec:nlp-extraction-merge}

对齐后的完全相同事实可按“身份、字段、值、条件、有效时间”规范键去重，
哈希仅用于索引，最终仍比较键内容并合并来源列表。
近重复记录先用受约束凝聚聚类确认身份，再合并互补字段；
并查集适合维护已经确认的等价关系，不适合对任意高相似边直接取闭包。
“同名且同地点”可能串起多周活动，角色值互补也不意味着属于同一实例。

冲突值不能总靠多数票解决。对于同一对象、同一字段、同一适用条件下的分类式候选值，
来源加权投票可写为
$\widehat{v}=\argmax_v\sum_jw_j\mathbb{I}[y_j=v]$。
权重应由来源核验、标注历史或开发集估计，不按报道篇数自动设为可靠性。
转载同一通知的十个页面并非十次独立观察，需要记录它们的依赖来源。
日期更正又是时间上的状态改变，应先拆开适用期，而不是把新旧排期交给投票。

当有许多同类记录和多个有误差的来源时，可用Dawid--Skene式模型同时估计潜在类别
与来源混淆矩阵\scite{dawid1979observer}
设记录$i$的潜在值为$z_i\in\{1,\ldots,K\}$，类别先验为$\pi_k$，
来源$j$的观测为$y_{i,j}$，并定义
$C^{(j)}_{k,\ell}=P(y_{i,j}=\ell\mid z_i=k)$。
在给定真实值后各来源条件独立的假设下，EM的E步为
\begin{equation}
  q_{i,k}=P(z_i=k\mid \bm{y}_i)=
  \frac{\pi_k\prod_{j\in O_i}C^{(j)}_{k,y_{i,j}}}
       {\sum_{k'}\pi_{k'}\prod_{j\in O_i}C^{(j)}_{k',y_{i,j}}},
  \label{eq:nlp-extraction-ds-e}
\end{equation}
$O_i$只含实际给出观测的来源；缺失不当作一个反对票。
M步用后验软计数更新
\begin{equation}
  \pi_k^{\mathrm{new}}=\frac{1}{N}\sum_iq_{i,k},
  \qquad
  C^{(j),\mathrm{new}}_{k,\ell}=
  \frac{\sum_{i:j\in O_i}q_{i,k}\mathbb{I}[y_{i,j}=\ell]}
       {\sum_{i:j\in O_i}q_{i,k}}.
  \label{eq:nlp-extraction-ds-m}
\end{equation}
实际可加Dirichlet平滑，避免少量数据把概率推到0或1；
按观测似然或参数变化判定收敛，多次初始化检查局部解。

\begin{example}[一次来源后验与参数更新]
\label{ex:nlp-extraction-ds}
构造两条静态二分类记录，值为A或B，不把可变排期放进这个静态真值模型。
初始先验各为0.5，两个来源的初始对称准确率分别为0.9和0.8。
记录1收到来源1的A与来源2的B，
A、B的未归一化权重为$0.5\times0.9\times0.2=0.09$与
$0.5\times0.1\times0.8=0.04$，所以$q_{1,A}=9/13\approx0.6923$。
记录2的两来源均报B，权重为$0.5\times0.1\times0.2=0.01$与
$0.5\times0.9\times0.8=0.36$，所以$q_{2,A}=1/37\approx0.0270$。

M步更新$\pi_A=(9/13+1/37)/2\approx0.3597$，
来源1在潜在A行报告A的概率变为
$(9/13)/(9/13+1/37)=333/346\approx0.9624$。
来源2只在这两条记录上报告过B，未经平滑的M步会把它每行报告B的概率都推为1。
这不是来源2已经被可靠辨认，而是样本太少、观测退化的提醒。
两条记录足以核对E/M运算，不足以可靠估计两个来源的全部误差参数。
\end{example}

没有锚定标签或足够的独立信息时，潜在类别可能发生标签置换，
高可信来源与一致但系统性错误的来源也可能难以区分。
转载违反条件独立假设，来源选择性报道又可能违反缺失机制假设。
因此后验是模型和数据条件下的置信判断，不是事实证明；
无法解释的冲突应并存保存并进入待核验状态。最新一条材料还须链接到正确活动，
不能仅因发布日期较晚便覆盖同名但不同实例的记录。

\subsection{记录核验与更新}
\label{sec:nlp-extraction-update}

每次写入前按模式、原文支持、对象身份、限定条件和来源完整性检查。
JSON Schema与SHACL复用第~\ref{sec:nlp-extraction-relation-records}节的方法，
字符偏移必须对应同一版本原文。支持模型判为蕴含后，
数字、否定和身份仍独立比较；这些局部错误常只差一个词或一位数字。
来源字段还应区分原始通知、转载、推断规则和人工确认，便于依赖回溯。

以图~\ref{fig:nlp-extraction-records}的教学通知为例：
4月14日旧材料将\code{ML-0418}排在4月17日，
4月15日更正为4月18日14:00。先按可靠活动编号确认是同一活动，
再解析“原定”与“现改至”指向同一个排期字段的旧值和新值。
输出不是删除旧来源，而是新增更正主张，连接其替代的旧主张，
同时保留改期事件及两份原文证据。原通知未给旧时刻，不应反填14:00。

\term{有效时间}（valid time）描述记录主张在所建模世界中适用的时间，
\term{写入时间}（transaction time）描述系统何时获知或写入这一版本。
在本例“更正自发布日生效”的明确约定下，
旧排期主张的有效期为$[\text{4月14日},\text{4月15日})$，
新排期为$[\text{4月15日},\text{尚未失效})$；
活动开始时刻4月18日14:00是排期字段的值，不是排期主张有效期的起点。
若系统到4月16日才读到更正，写入日期为4月16日，
仍可记录更正自4月15日起有效。另有一份4月17日才导入的旧通知，
不能因为写入较晚就恢复旧日期。

增量差分按“身份、字段、条件、适用期”比较新旧记录。
缺少的新字段是补充，同字段同条件的新值可能是修改，
明确否认旧主张的是撤回，预设适用期结束是过期；
当前材料未提某字段并不等于撤回该字段。
更正只涉及活动日期时，报名截止仍为4月16日18:00前，
不能按日期差自动顺延。若更正对截止条件也有说明，则单独更新并记录相应证据。

来源撤回时，沿“来源版本$\to$抽取主张$\to$合并记录$\to$派生关系”的依赖边定位受影响项。
若另有独立来源支持，不必立即删除事实，但要重算支持状态；
仅由撤回来源支撑的派生记录则需降级或撤回，保留操作历史。
这与直接删除整份文档中的全部对象不同，实体存在与某条关系成立是两种主张。

评价应同时报告事实正确性、覆盖、重复率、未决冲突和证据位置正确性，
再按抽取、绑定、归一化、归并、更新定位误差。
漏掉报名截止会使问答系统给出错误参加条件，合并错两场活动则会返回别场地点，
即使最终回答语言通顺也无法弥补记录错误。
版本索引、事务和服务性能属于系统实现问题，本章确定的是这些实现必须保留的
身份、条件、来源及时间语义。

\section*{文献说明}

Few-NERD、DocRED、MAVEN、MAVEN-ERE在本章用于说明任务与监督对象；
W2NER、UIE、CasRel、TPLinker、EIDER、DMCNN、Text2Event和EEQA用于说明具体预测与恢复机制。
以下给出这些资料的完整出处。
\begin{fullwidth}
  \small
  \noindent\fullcite{ding2021fewnerd}\par\medskip
  \noindent\fullcite{li2022w2ner}\par\medskip
  \noindent\fullcite{lu2022uie}\par\medskip
  \noindent\fullcite{yao2019docred}\par\medskip
  \noindent\fullcite{wei2020casrel}\par\medskip
  \noindent\fullcite{wang2020tplinker}\par\medskip
  \noindent\fullcite{xie2022eider}\par\medskip
  \noindent\fullcite{wang2020maven}\par\medskip
  \noindent\fullcite{chen2015dmcnn}\par\medskip
  \noindent\fullcite{lu2021text2event}\par\medskip
  \noindent\fullcite{du2020eeqa}\par\medskip
  \noindent\fullcite{wang2022mavenere}
\end{fullwidth}

\section{本章小结}
\label{sec:nlp-extraction-summary}

从一段通知到可查询记录，关键不在于输出了多少名字或三元组，
而在于每一项输出能否定位、绑定和核验。
提及提供原文位置，共指与链接分别确定文档内同指关系和知识库身份；
关系把对象与表达绑定，事件把角色与同一次发生或计划绑定，
时间、因果和组成关系再连接不同事件。丢掉任何一层区分，
都可能得到格式正确却指向错误对象、日期或状态的结果。

序列、跨度、词元对和生成模型提供不同预测结构，
候选裁剪、模式约束及全局优化帮助恢复合法结果，
却不能取代原文对事实的支持。证据选择需要独立监督与检查，
阈值和来源模型也只在相应数据假设下提供判断。
可用的知识记录因此还必须保存条件、来源和有效时间，
在新材料到达时区分补充、更正与撤回，而不是简单追加或覆盖。

章首通知中的4月17日应作为旧排期保留，4月18日14:00是当前活动开始时刻，
4月16日18:00前则限定校外报名。三项日期并不竞争同一个字段，
只有先明确对象、角色与状态，才能回答何时举行、如何参加，以及答案依据是什么。
下一章转向文档与文本集合的分类、匹配和组织，
本章形成的实体、关系和带来源记录可以作为其中的结构化输入，
其不确定性与适用边界也应一并传递。
